\section{F. The later Montpellier topoi}
\label{sec:F}

We distinguish three constructions throughout this section: a pro-object of
full neighbourhood topoi, the classifying topos of covering germs, and a topos
of sheaves constructible for a fixed stratification. They have different
universal properties. We write $B G$ for the topos of continuous discrete
sets with an action of a profinite group $G$, and $B\mathcal G$ for its
groupoid version. A morphism over a base topos includes its commutation
\emph{isomorphism}; forgetting that isomorphism changes several answers below.

\subsection{F(a): the six presentations and their rigidifications}

\paragraph{F.1 (Proved). The six presentations, with the information retained.}
Fix a profinite group $\Gamma$ and its one-object profinite groupoid
$B_{\mathrm{gr}}\Gamma$. The following are presentations of the same
unpointed object, provided the indicated two-dimensional data are retained.
\begin{enumerate}
\item Connected point-admitting profinite Galois topoi with a geometric morphism to
$B\Gamma$ whose induced group homomorphism has open image. No point is
selected. Morphisms are geometric morphisms together with a specified
isomorphism over $B\Gamma$; two-morphisms respect that isomorphism.
\item Connected profinite groupoids over $B_{\mathrm{gr}}\Gamma$ with open image, with functors,
commutation isomorphisms, and compatible natural isomorphisms.
\item One-object presentations of (2), modulo inner conjugacy, \emph{together
with the conjugating elements as two-morphisms}. If only exterior classes
are retained, this is the homotopy category, not the original two-category.
\item Homomorphisms $p:E\to\Gamma$ of profinite groups with open image. A morphism from
$p':E'\to\Gamma$ to $p:E\to\Gamma$ is a pair $(u,\gamma)$ with
$u:E'\to E$ continuous and $p u=\operatorname{Ad}(\gamma)p'$.
A two-morphism $e:(u,\gamma)\Rightarrow(v,\delta)$ satisfies
$v=\operatorname{Ad}(e)u$ and $\delta=p(e)\gamma$ (the equivalent inverse
convention gives the same two-category).
\item If $H=p(E)$ is open, the extension $1\to\pi\to E\to H\to1$, together
with its embedding $H\hookrightarrow\Gamma$. If $Z(\pi)=1$, one can instead
retain the continuous outer action $H\to\operatorname{Out}(\pi)$, with its
extension topology; morphisms require the intertwiners and their coherence.
\item A profinite groupoid with continuous weak $\Gamma$-action, whose
homotopy quotient is connected and whose object-component $\Gamma$-set is
finite; equivariant functors and equivariant natural isomorphisms are its
morphisms. Strict actions may be used after the usual transport-groupoid
strictification, without discarding its coherence.
\end{enumerate}
Without the open-image restriction, replace the finite component set in (6)
by an arbitrary profinite homogeneous $\Gamma$-space, and omit the word \emph{open} in (5).
These descriptions are equivalent for the data arising from a finitely
generated field $K/\mathbb Q$. In (5), $H=\Gamma_k$, where $k$ is the
algebraic closure of $\mathbb Q$ in $K$, embedded in a fixed
$\overline{\mathbb Q}$; in (6) the groupoid is that of
$K\otimes_{\mathbb Q}\overline{\mathbb Q}$. If $K$ is finite over
$\mathbb Q$, then $\pi=1$: the subgroup $H$, or equivalently the finite
$\Gamma$-set $\Gamma/H$, must still be retained. The trivial geometric group
alone does not distinguish number fields.

\begin{proof}
For a profinite groupoid, evaluation at each object is a fibre functor on
continuous actions. Equivariant natural transformations of evaluation are
exactly the continuous arrows of the groupoid: finite transitive actions
detect each finite quotient, and their inverse limit detects the arrow.
Conversely the fibre functors of a pointed connected profinite Galois topos
recover the finite Galois category and its continuous-action topos. This
proves (1)$\simeq$(2), including all natural isomorphisms, without choosing
one object globally. Choosing an object of a connected groupoid gives its
one-object skeleton; changing that object is an isomorphism torsor, which
is exactly the inner-conjugacy datum in (3)--(4). On one-object groupoids
the naturality equation is the displayed conjugacy equation; compatibility
with the base is its second displayed equation. Consequently forgetting
the conjugating elements loses the centralizer of $u(E')$ intersected with
$\ker p$, and is a biequivalence only on Hom-groupoids where that group is
trivial.

The geometric fibre of $B_{\mathrm{gr}}E\to B_{\mathrm{gr}}\Gamma$ is the
transport groupoid with objects $\Gamma$ and arrows given by right
translation by $E$ through $p$. Its components are $\Gamma/H$, its vertex
groups are $\pi$, and left translation supplies the $\Gamma$-action.
Taking the transport homotopy quotient reverses this construction. This is
also a direct verification of effective Galois descent of actions, so it
proves (4)$\simeq$(6), rather than identifying a quotient group with a
quotient groupoid.

For a centreless kernel, conjugation identifies an extension with
\[
 E\ \simeq\ \operatorname{Aut}(\pi)
       \mathbin{\times}_{\operatorname{Out}(\pi)} H.
\]
Indeed its kernel is $Z(\pi)$, and any pair on the right differs from the
conjugation action of a lift of its $H$-coordinate by a unique inner
automorphism. This proves both injectivity and surjectivity and also the
uniqueness of the extension isomorphism. In the profinite situation all
maps are continuous; the extension topology is retained, so the assertion
does not silently replace continuous extensions by abstract extensions.
For a general kernel the outer action alone has the usual central
extension ambiguity: already $H=C_2$, $\pi=C_2$ and trivial outer action
admit $C_4$ and $C_2\times C_2$. This proves exactly what description (5)
would otherwise lose. Finally the components of
$K\otimes_{\mathbb Q}\overline{\mathbb Q}$ are indexed by embeddings
$k\hookrightarrow\overline{\mathbb Q}$, and their stabilizer is $\Gamma_k$.
\end{proof}

\paragraph{F.2 (Proved). Openness and the centralizer used at scan 20.}
Let $F\subset F'$ be finitely generated extensions of an algebraically
closed characteristic-zero field. Then the image of
$G_{F'}\to G_F$ is open. Every open subgroup of $G_F$ has trivial
centralizer in $G_F$ when $\operatorname{trdeg}F>0$. For transcendence
degree zero, $G_F=1$ and its centralizer is again the trivial group.
Consequently the geometric exterior homomorphism at MPL~149, fol.~9,
scan~20 determines a geometric morphism of classifying topoi up to unique
isomorphism. This assertion is about the geometric groups; openness alone
is not a centralizer theorem for arbitrary profinite groups.

\begin{proof}
The relative algebraic closure $L$ of $F$ in $F'$ is finite over $F$.
The regular extension $F'/L$ is linearly disjoint from every finite
separable extension of $L$. Every automorphism of $L^{\mathrm{sep}}/L$
therefore extends to the compositum with $F'$, and then to a separable
closure of $F'$. Thus the image is $G_L$, an open subgroup of $G_F$.

We give the needed slimness argument. For every finite Galois extension
$L/F$ with group $Q$ and every prime $\ell$, there exists a finite Galois
extension $M/F$ above it such that the action of $Q$ on
$N=\operatorname{Gal}(M/L)$ is faithful. Choose a normal model of $F$ and
shrink it so that the normalization in $L$ is finite etale. A closed point
in its smooth locus has a full $Q$-orbit above it, since the constant field
is algebraically closed. A divisorial valuation centred there (use the
exceptional divisor of the blow-up when the dimension exceeds one) has
$|Q|$ distinct prolongations to $L$. Approximation for these finitely many
independent discrete valuations gives $a\in L^\times$ with valuation one
at one prolongation and zero at all the others. Its conjugates are
independent in $L^\times/L^{\times\ell}$, as their valuation vectors are
the standard basis. Since $\mu_\ell\subset F$, adjoining all their
$\ell$th roots gives $N\simeq(C_\ell)^Q$, with faithful permutation action.
If $z$ is central in $G_F$, its image in $Q$ acts trivially on $N$,
so faithfulness gives that image equal to one. Since $L/F$ was arbitrary,
$z=1$. The same argument applies to every open subgroup, because its fixed
field is again finitely generated over the same algebraically closed
constant field. Finally, if $z$ centralizes an open subgroup $H$, then
$J=\langle H,z\rangle$ is open, $H$ is normal in $J$, and $z$ lies in
$Z(J)$. The already proved centrelessness of $J$ gives $z=1$. Thus every
open subgroup has trivial centralizer in $G_F$. Natural automorphisms of a restriction functor are precisely
the centralizer of the image, proving the final assertion.
\end{proof}

\paragraph{F.3 (Proved). Faithfulness of fields in every characteristic.}
Let $k$ be any field and $K,K'$ finitely generated extensions. Two
$k$-embeddings $f,g:K\to K'$ inducing isomorphic morphisms of profinite
Galois groupoids \emph{over} $\Pi_k$ are equal. Thus the characteristic-zero
margin to Theorem~2, MPL~149, fol.~10, scan~22, is correct; a
positive-characteristic counterexample with the same relative data does
not exist.

\begin{proof}
For each integer $n$ invertible in $k$, the coefficient $\mu_n$ is pulled
back from $B_k$. The specified relative isomorphism therefore identifies
the pullbacks in Kummer cohomology, with the specified identification of
coefficients, so
\[
 f(a)/g(a)\in K'^{\times n}\qquad(a\in K^\times).
\]
The same remains true after any algebraic extension of the constants and
on each connected component of the geometric base change. Choose a
component over an algebraic closure $\bar k$ of $k$ and write $F'$ for its
function field. For $a$ transcendental over $k$, both $f(a)$ and $g(a)$
remain nonconstant in $F'$. On a normal proper integral $\bar k$-model of
$F'$, every prime-divisor valuation of their ratio is divisible by all
integers prime to the characteristic, and hence zero. A rational function
with zero divisor on a normal proper integral variety over $\bar k$ is a
constant: regularity in codimension one and normality make it an everywhere
regular unit, and properness makes it an element of $\bar k^\times$.
Consequently $g(a)=\lambda f(a)$. Applying the same argument to $a-1$
gives $g(a)-1=\mu(f(a)-1)$. Nonconstancy gives $\lambda=\mu=1$.

For $a$ algebraic over $k$, the finite separable part of $k(a)/k$ is an
object of the finite etale category of $k$. The relative commutation datum
recovers its embedding in $K'$ by the universal property of the associated
slice topos. A purely inseparable extension has at most one embedding
extending that separable embedding. Thus $f(a)=g(a)$ also in this case.
The map from $K'$ into the chosen geometric component field is injective
(it is a map from a field), so equality there implies equality in $K'$.
\end{proof}

\paragraph{F.4 (Proved). Relative faithfulness of reduced schemes.}
Let $S$ be locally noetherian and let $X,Y$ be locally of finite
presentation over $S$, with $X$ reduced. If $f,g:X\to Y$ induce isomorphic
morphisms $X_{\mathrm{et}}\to Y_{\mathrm{et}}$ in the two-category over
$S_{\mathrm{et}}$, then $f=g$. No characteristic-zero hypothesis and no
reducedness hypothesis on $Y$ are necessary. In particular the requested
positive-characteristic refutation of this relative statement cannot be
given.

\begin{proof}
First suppose $S=\operatorname{Spec}k$. Isomorphic geometric morphisms
induce the same map on subterminal objects and hence the same map on
underlying Zariski spaces. At a closed point $x\in X$ their common image
$y\in Y$ is closed, and both residue fields are finite over $k$. The two
maps on residue fields have the same effect on the maximal separable
subextension of $\kappa(y)/k$: a finite separable extension represents a
finite etale object of $B_k$, and the slice-topos universal property
recovers its actual map, not merely a conjugacy class. The remaining
purely inseparable extension has a unique embedding. Thus $f$ and $g$
agree on every closed residue point. Locally on a common inverse image of
an affine open in $Y$, their difference on any coordinate function vanishes
in every closed residue field. The source is reduced and Jacobson, so the
intersection of its maximal ideals is zero. Hence the two ring maps, and
therefore the morphisms, agree.

For general $S$, the square of etale topoi obtained by base change along
$\operatorname{Spec}\kappa(s)\to S$ is cartesian. One can verify this by
writing the local scheme as the inverse limit of affine open
neighbourhoods of $s$ and then taking its closed residue subtopos; this is
also \cite[Lemma~2.9]{CHW2024}. Pull the relative isomorphism to this square,
and then restrict to $(X_s)_{\mathrm{red}}$. The field case applies to this
reduced finite-type source and proves equality on every reduced fibre.
The morphism
$\coprod_s(X_s)_{\mathrm{red}}\to X$ is schematically dominant for reduced
$X$: a function vanishing on every such fibre vanishes in every residue
field and hence is zero. Since the underlying point maps already agree,
the same affine-local argument proves $f=g$. All reductions are local on
$S,X,Y$, so no global quasi-compactness is required.

For comparison with \cite[Theorem~3.9]{CHW2024}, the extra reduced-fibre
hypothesis in its stated version can be removed here by explicitly
restricting to the reduced fibres, as above. This argument asserts
faithfulness, not fullness.
\end{proof}

\paragraph{F.5 (Refuted). Unrestricted
full faithfulness is false; pinned reconstruction is the correct theorem.}
For every characteristic-zero field $k$, morphisms of etale topoi over
$B_k$ need not arise from scheme morphisms, even after universal
homeomorphisms are inverted. If $k$ is an infinite finitely generated
field, the functor becomes fully faithful into the category of
\emph{pinned} geometric morphisms, which send closed points to closed
points. The same holds for sub-$p$-adic characteristic-zero fields.

\begin{proof}
The $t$-adic decomposition subgroup gives a section of
$G_{k(t)}\to G_k$: explicitly embed $k(t)$ in $k((t))$, use the splitting
of $G_{k((t))}\to G_k$ obtained by acting on coefficients and fixing
compatible roots of $t$, and restrict to a separable closure of $k(t)$.
The resulting morphism $B_k\to B_{k(t)}\to(\mathbb P^1_k)_{\mathrm{et}}$
over $B_k$ has image the generic point. A $k$-point of $\mathbb P^1_k$
has closed image, and universal homeomorphisms do not change that
obstruction. This is the explicit obstruction in
\cite[Example~2.30]{CHW2024}.

For an infinite field finitely generated over its prime field,
\cite[Theorem~0.1]{CHW2024} gives an equivalence on the Hom-sets of the
localization at universal homeomorphisms and the pinned Hom-groupoids.
The latter are equivalent to sets by its Proposition~2.22. The source
localization is the category of absolute weak normalizations: in
characteristic zero this is seminormalization, and in positive
characteristic it is perfection. Thus the theorem applies to the
manuscript's reduced finite-type schemes with exactly this localization
and this target restriction. The assertion for sub-$p$-adic fields is
\cite[Theorem~A, equivalently Theorem~4.1]{CarlsonStix2024}; its
identification of the localization with seminormal schemes gives the
same reduction. These statements do not classify arbitrary base fields.
\end{proof}

\paragraph{F.6 (Refuted). Even pinned reconstruction fails over every
algebraically closed characteristic-zero field.}
Let $k$ be algebraically closed of characteristic zero. Put $j=i$ and
$j'=-i$, and let
\[
 E_j:\quad y^2=x^3-\frac{3j}{j-1728}x-\frac{2j}{j-1728}.
\]
Then $E_j$ and $E_{j'}$ are nonisomorphic smooth proper $k$-curves but
have equivalent etale topoi over $B_k\simeq\mathbf{Set}$. This remains a
counterexample after universal homeomorphisms are inverted.

\begin{proof}
The displayed Weierstrass equation has nonzero discriminant and
$j$-invariant $j$; the symbolic computation is in
\texttt{code/montpellier\_checks.py}. Extend the automorphism
$i\mapsto-i$ of $\overline{\mathbb Q}$ to $k$, by choosing and fixing a
transcendence basis and then extending to an algebraic closure. Applying
this automorphism to coefficients gives an isomorphism of the underlying
schemes and hence an equivalence of their small etale topoi. It is over
$\mathbf{Set}$ automatically and is pinned because it is a homeomorphism.
A $k$-isomorphism of the curves would preserve the $j$-invariant,
contradicting $i\ne-i$. The curves are normal, hence seminormal, so their
objects in the localized characteristic-zero category remain themselves.
\end{proof}

\paragraph{F.7 (Known). The experiment with a proper
curve over a number field.}
Let $k$ be a number field, $C/k$ a smooth proper geometrically connected
curve of any genus, $K=k(C)$, and
$\pi=G_{K\bar k}$. The $k$-automorphisms of $C$ are in bijection with
continuous outer automorphisms of $G_K$ over $G_k$, and equivalently with
continuous outer automorphisms of $\pi$ commuting with the outer
$G_k$-action. Requiring preservation of the conjugacy classes of the
inertia groups of closed points does not change this answer. Smoothness
and normality are essential if the curve, rather than its normalization,
is to be recovered from its function field.

\begin{proof}
Apply \cite[Theorem~B, Theorem~17.1 in the body]{Mochizuki1999} to the
regular function field $K/k$ on both sides; $k$ is sub-$p$-adic. An
isomorphism of groups is open, so it and its inverse give field embeddings.
Their composites are the identity by the theorem's uniqueness, hence
are inverse $k$-automorphisms. A birational map between smooth proper
curves extends uniquely at every discrete valuation ring by the valuative
criterion of properness, and the inverse extends as well. Thus it is an
automorphism of $C$. Conversely any such automorphism induces the stated
outer group automorphism and permutes the local valuations and their
inertia. F.2 makes $\pi$ centreless; the extension pullback in F.1 therefore
turns a $G_k$-equivariant outer automorphism of $\pi$ into the required
outer automorphism of $G_K$. This verifies the proposed experiment at
MPL~149, \S2, fol.~25v, scan~52, including genus zero and genus one; it uses the
birational group, not just the proper curve's fundamental group.
\end{proof}

\subsection{F(c): quasi-Galois gluing and ordinary exit categories}

\paragraph{F.8 (Proved). A precise finite quasi-Galois characterization.}
Suppose a Grothendieck topos $\mathcal T$ has a finite filtration by open
subtopoi, each successive locally closed stratum being the presheaf topos
of a groupoid with finitely many connected components. At each step write
$F=i^*j_*:\mathcal T_{\rm open}\to\mathcal T_{\rm closed}$ for the
accessible left-exact gluing functor. Then $\mathcal T$ is, compatibly with
the filtration, a presheaf topos on an EI-category with finitely many
isomorphism classes if and only if every $F$ preserves arbitrary products.
Here an EI-category is a category in which every endomorphism is
invertible, and its isomorphism classes are partially ordered by the
existence of a morphism. Without the finiteness of the component sets of
the strata, the corresponding conclusion uses a small EI-category and
need not have finitely many isomorphism classes.

\begin{proof}
We prove one induction step. Write the open topos as
$\widehat A=\operatorname{Fun}(A^{\rm op},\mathbf{Set})$ and the closed
one as $\widehat B$, with $B$ a groupoid. A sheaf in the glued topos is a
triple $(P,Q,q:Q\to F(P))$. If $F$ preserves products, left exactness
makes it preserve every small limit. Accessibility and the adjoint
functor theorem give a left adjoint $L$ to $F$. For each representable
$h_b\in\widehat B$ put $M_b=L(h_b)\in\widehat A$. Yoneda and adjunction
give
\[
 F(P)(b)=\operatorname{Hom}_{\widehat B}(h_b,F(P))
          \simeq\operatorname{Hom}_{\widehat A}(M_b,P).
\]
Define a category $C$ by adjoining to $A$ the objects of $B$, putting
$\operatorname{Hom}_C(a,b)=M_b(a)$, no morphisms from $B$ to $A$, and
using the actions supplied by the presheaves $M_b$ and functoriality in
$b$ for composition. A presheaf on $C$ is exactly the triple above.
This proves $\mathcal T\simeq\widehat C$. All new endomorphisms lie in
$A$ or $B$, so the EI property follows by induction. No cycle can cross
the finite filtration; after passing to isomorphism classes the resulting
preorder is a partial order.

Conversely, for an EI-category with the given filtration, crossing arrows
form the bimodule $M_b(a)=\operatorname{Hom}_C(a,b)$. The gluing functor
has the displayed representable form and preserves all limits,
particularly products. The numbers of isomorphism classes add under the
finite construction. A single stratum $\mathbf{Set}^{\mathbb N}$ shows
why that finite-component condition cannot be omitted: its indecomposable
projective generators have infinitely many nonisomorphic types.
This proves the categorical case asserted in STRATA B.1--B.2 rather than
assuming that arbitrary left-exact gluing is representable by crossing
sets.
\end{proof}

\paragraph{F.9 (Proved). A two-stratum gluing which is not of presheaf EI type.}
Let $C=\{0,1\}^{\mathbb N}$ be the Cantor space, and define
$F(A)=\operatorname{LC}(C,A)$ for a set $A$ with the discrete topology.
The Artin gluing of two copies of $\mathbf{Set}$ along $F$ is a
Grothendieck topos with two Galois strata, but it is not the presheaf
EI-topos asserted in F.8.

\begin{proof}
Every locally constant map from compact $C$ has finite image and factors
through a finite clopen partition. Consequently
$F=\operatorname*{colim}_{\mathcal P}(-)^{\mathcal P}$, where
$\mathcal P$ runs through the filtered set of finite clopen partitions.
It is accessible and preserves finite limits, so its Artin gluing is a
Grothendieck topos. For the countable family $A_n=\{0,1\}$, the coordinate
projections $C\to A_n$ are locally constant. They define an element of
$\prod_nF(A_n)$ which does not come from $F(\prod_n A_n)$: the only map
with these coordinates is the identity of the underlying set of $C$,
which is not continuous from Cantor space to that set with the discrete
topology. Thus $F$ fails to preserve this product, and F.8 proves the
claim. This example also specifies the missing gluing condition in the
undefined term ``topos quasi-galoisien'' at MPL~147, fol.~85, scan~179.
\end{proof}

\paragraph{F.10 (Proved). Constructible sets and the crossing double cosets.}
Let $(X,D)$ be a complex manifold with a locally finite normal-crossings
divisor, or an analytic Deligne--Mumford stack with such an etale local
presentation. Assume the usual paracompact local-conical hypotheses for
constructible exodromy. Write $\mathcal E=\operatorname{Exit}(X,D)$,
with arrows from closed to more open strata. Then
\[
 \operatorname{Sh}_{\rm str}(X;\mathbf{Set})
       \simeq\operatorname{Fun}(h\mathcal E,\mathbf{Set})
       =\widehat{(h\mathcal E)^{\rm op}}.
\]
This is a quasi-Galois topos in the sense of F.8 (without finite component
restriction if necessary). For one connected incidence link $L$ between
connected strata $S,T$, choose endpoints and put
$A=\pi_1S$, $B=\pi_1T$, $G=\pi_1L$. The crossing-arrow set, with its
$A$--$B$ actions, is
\[
             (A\times B)/\operatorname{im}(G\to A\times B),
\]
with the harmless opposite convention on one factor dictated by endpoint
transport. If several link components occur, take the disjoint union of
these orbit sets. Composition is induced by concatenation on the full
flag spaces, not by composing their sets of components without the flag
maps.

\begin{proof}
A constructible sheaf transports germs along paths within each stratum.
In a conical chart a generization from a closed stratum is transport along
an outgoing radial arc; homotopic exit arcs induce the same map of germs.
The sheaf axiom on conical charts and their intersections gives precisely
compatibility under concatenation, and produces a functor on
$h\mathcal E$. Conversely, for such a functor assign on a conical open
chart the compatible germs indexed by its exit objects. Restriction is
restriction of compatible families. The charts form a basis stable under
refinement, and compatibility on a cover is exactly the sheaf condition;
the stalk at an exit object is the assigned set. On each stratum the
construction is the usual locally constant sheaf associated to a
fundamental-groupoid action. These two constructions are inverse on
stalks, hence on sheaves, and on natural transformations. For the stack
apply this construction on the entire etale atlas nerve and descend:
compatible charts include the stabilizer actions. This is the
set-valued reduction of the exodromy theorem, with all transition data
retained.

The exit mapping space for this incidence is the homotopy fibre of
$L\to S\times T$ at the chosen endpoints. The long exact sequence in
components of that fibre gives the orbit-set formula. Its residual
endpoint actions give the two module actions. Two composable exits are
represented by a flag with two steps; the face forgetting the middle
step is their composition. Three-step flags verify associativity. Within
a stratum every exit is invertible, and a strict exit cannot return to its
source stratum, so all endomorphisms are invertible. Finally the bimodule
formula for restrictions shows that its gluing functor is represented as
in F.8. This verifies the relevant STRATA C.1 and C.6 calculation in the
geometric setting; it does not identify the full sheaf topos of $X$ with
this constructible subcategory.
\end{proof}

\paragraph{F.11 (Refuted). The Boolean source axiom does not imply thinness.}
The weighted Boolean axiom (174)c at MPL~147, fols.~56--57,
scans~115--116, implies (174)a and b. It does \emph{not} imply that an
arrow is determined by its source and target. The additional geometric
transversality hypothesis later imposed at scans~120--121 gives this
property in the intended branch-incidence examples.

\begin{proof}
The coslice of a rank-$d$ object has the weighted Boolean lattice of
subsets of a $d$-element set. Its ranks are between zero and $d$, and it
has a unique object of rank $d$, namely the identity. This proves a and b.
For the false extra implication use the finite category with objects
$x,y,z$ of ranks $2,1,0$, arrows
\[
 a,b:x\longrightarrow y,\qquad c:y\longrightarrow z,
 \qquad d:x\longrightarrow z,
 \qquad ca=cb=d,
\]
and identities. Its coslices are the Boolean diamond, the two-element
chain, and the one-point poset, respectively. Thus it satisfies the full
weighted Boolean coslice condition, but $a\ne b$ have identical source
and target. Regard all these finite sets as discrete zero-dimensional
analytic spaces to obtain an internal analytic category. The program
\texttt{montpellier\_checks.py} verifies every composition and every
coslice relation.

The exact correction to the axiom is injectivity of the target map on
every source coslice, or equivalently monicity of
$(s,t):D_{**}\to D_*\times D_*$. For genuine normalized branch germs, an
object is a specified set of local branches, so an arrow exists exactly
when one specified set contains the other, and its inclusion is unique.
Under the later transverse geometric realization, distinct subsets give
distinct intersection germs by independent local coordinates, so this
extra condition holds. It is the geometric realization, not the
abstract Boolean axiom alone, that supplies it.
\end{proof}

\subsection{F(d): germs, full sheaves, and descent}

\paragraph{F.12 (Proved). The punctured germ is a pro-topos; its covering
reflection is a different object.}
For a locally closed analytic embedding $Z\hookrightarrow Y$, let
$\mathfrak V_{Z,Y}$ be the pro-object of full sheaf topoi
$(\operatorname{Sh}(V))_V$ indexed by open neighbourhoods of $Z$ in an
ambient open where $Z$ is closed. With a divisor $D$, put
$\mathfrak V^*_{Z,Y}=(\operatorname{Sh}(V\setminus D))_V$.
These definitions are invariant under cofinal refinement. Their
covering-germ categories are the filtered two-colimits of the covering
categories at the stages. For a normal-crossing point with $r$ branches
this category is $\mathbb Z^r$-sets, and its finite part is finite
continuous $\widehat{\mathbb Z}(1)^r$-sets. In contrast the literal
inverse two-limit of the full sheaf topoi of shrinking punctured discs
is degenerate.

\begin{proof}
Cofinal systems present the same pro-object by the definition of
pro-morphisms. Polar coordinates identify a cofinal system of punctured
normal polydiscs with a contractible radial factor times $(S^1)^r$;
shrinking acts identically on the coordinate winding numbers. Covering
classification therefore gives the asserted filtered category, including
its morphisms after common restriction. This verifies all steps of
LM8 E.10 independently.

For the last assertion let $D_n^*=\{0<|z|<1/n\}$ and suppose a topos
$\mathcal A$ maps compatibly to every $\operatorname{Sh}(D_n^*)$. Viewed
as opens of $D_1^*$, the annuli $U_k=\{1/k<|z|<1\}$ cover $D_1^*$.
Each inverse image of $U_k$ is initial because the morphism factors
through $D_k^*$, disjoint from $U_k$. Inverse images preserve unions,
so the terminal object of $\mathcal A$ is initial. The degenerate topos
has the unique such cone, proving the universal property. This also
verifies LM8 E.11. The manuscript's arcs and $C^\infty$ embeddings at
scans~143--148 identify these as classical analytic germs; a profinite
etale replacement is a separate construction, not a new reading of the
word ``germe''.
\end{proof}

\paragraph{F.13 (Proved). The precise full-topos gluing theorem.}
Let $X$ be paracompact and locally compact, $D$ a locally finite
normal-crossings divisor, and $\Sigma\subset D$ a closed union of
strata. A cofinal system of controlled tubular neighbourhoods of $\Sigma$
can be covered by tubular neighbourhoods of its normalized branches,
with all their intersections. The resulting full Cech diagram, or its
full flag subdivision, has two-colimit
$\mathfrak V^*_{\Sigma,X}$ in pro-Grothendieck topoi. This is the correct
full-sheaf meaning of both versions of the gluing theorem at MPL~147,
fols.~74 and 76, scans~150 and 155. The same theorem holds for analytic
Deligne--Mumford stacks presented by proper etale groupoids with
paracompact locally compact orbit space, or for an explicitly controlled
atlas whose descent conditions admit the common-stage refinement used
below. Arbitrary nonproper atlases, even with finite isotropy, require
that extra hypothesis: F.51 refutes passage from objectwise atlas germs to
an intrinsic stack neighbourhood without it. For a closed $D$ one also has
\[
 \operatorname{Sh}(X)\ \simeq\
 \operatorname{Sh}(X\setminus D)
       \amalg_{\mathfrak V^*_{D,X}}\mathfrak V_{D,X}
 \quad\text{in pro-topoi},
\]
where the left side is constant. This proves the full-topos version of
(254) and (266), with their germ objects correctly typed.

\begin{proof}
At a fixed neighbourhood stage $V$, choose open tubular sets $V_a$ whose
union is $V$ (and intersect them with $X\setminus D$ for the punctured
version). Sheaves on $V$ are exactly sheaves on the $V_a$ with
isomorphisms on double intersections satisfying the triple cocycle
condition. Here is the topoi-level universal property: for any target
$\mathcal A$, a geometric morphism $\operatorname{Sh}(V)\to\mathcal A$
is a finite-limit- and colimit-preserving functor
$\mathcal A\to\operatorname{Sh}(V)$. Restricting that functor produces
such functors on every $V_a$ with compatible isomorphisms. Conversely
objectwise sheaf descent glues these functors and natural transformations;
finite limits and colimits are checked after the conservative open
restriction family. Thus the gluing is a bicategorical colimit of
\emph{topoi}, not merely a statement about covering categories.

For finitely many local branches, subdividing the Cech index category by
all nonempty nested intersections does not change descent data: a
subdivided datum is recovered by restriction from its one-step data, and
the flag triangles are precisely the cocycle relations. For completeness the passage from local germs to a global pro-object
requires a uniformization argument. Fix a target topos $\mathcal A$ and a
compatible germ datum. On a locally finite, relatively compact shrinking
of the coordinate cover choose a representative neighbourhood for every
chart datum, every overlap isomorphism, and every triple cocycle equality.
Near any given point only finitely many of these conditions occur.
After shrinking the coordinate sets once more, their allowed radii have
a positive lower bound on each compact member of that shrinking. A
partition of unity gives a positive continuous function smaller than all
these locally finite bounds. Restricting to that variable-radius tube
realizes every datum and every compatibility at one stage. The same
argument applied to two representatives, and then to their two-morphisms,
proves full faithfulness as well as essential surjectivity. Thus finite
chart descent passes to germs and global open descent is legitimate;
we have not commuted a filtered colimit with an arbitrary infinite
limit. For a pro-target first apply this argument to each fixed target
stage, and then take the limit over target stages in the pro-mapping
formula. This proves the first assertion and its version for $\Sigma$.

For the second assertion, at every stage $V$ the pair
$X\setminus D,V$ is an open cover of $X$, with intersection
$V\setminus D$. Apply the just-proved two-open universal property and
pass through the same cofinal system. For non-simple normal crossings,
choose an etale chart atlas and all its fibre products; the normalized
local branches are handled on this atlas, not assumed globally embedded.
For a proper etale groupoid the local orbifold charts have finite
stabilizers. Choose a locally finite shrinking on the orbit space and
take the minimum over the finitely many stabilizer translates of each
chart radius. Over compact chart closures, properness makes the relevant
transition conditions finite after an etale refinement. The preceding
uniformization therefore applies to the overlap maps and their cocycles
as well as to the branch data. Tubular constructions and intersections
pull back etale locally, and these synchronized descent diagrams may be
composed. The result retains stabilizer groups and requires no invariant
tubular embedding. For a nonproper groupoid this finite-control step
must be supplied separately; F.51 gives a case where it fails. The
intrinsic pro-object indexed by actual stack neighbourhoods still makes
sense, and its two-open pushout is proved at each actual common stage;
it is not thereby identified with the objectwise atlas-germ diagram.
\end{proof}

\paragraph{F.14 (Refuted). Covering topoi cannot replace full elementary
topoi in an ordinary two-colimit.}
For discrete groupoids $\mathcal G_i$, the ordinary bicategorical colimit
of the topoi $B\mathcal G_i$ is
$B(\operatorname*{2\!colim}_i\mathcal G_i)$. It generally has the wrong
higher shape for the geometric gluing. In particular
\[
 \mathbf{Set}\amalg_{B\mathbb Z}\mathbf{Set}\simeq\mathbf{Set},
 \qquad *\amalg_{S^1}^{h}*\simeq S^2.
\]
Thus neither ``aspherical strata'' nor ``all flags'' repairs the
replacement of full elementary topoi by covering topoi inside the
ordinary two-category. One must instead take the homotopy colimit of
classifying \emph{spaces}, or use the corresponding infinity-topoi.

\begin{proof}
Present a topos $\mathcal A$ by a small site. A geometric morphism
$B\mathcal G\to\mathcal A$ is a flat continuous site functor valued in
$\mathcal G$-sets. Evaluation at the objects of $\mathcal G$ identifies
this with a $\mathcal G$-diagram of points of $\mathcal A$, with natural
transformations: flatness and covering epimorphisms are checked
objectwise. Hence
$\operatorname{Geom}(B\mathcal G,\mathcal A)
\simeq\operatorname{Fun}(\mathcal G,\operatorname{Pt}\mathcal A)$.
The bicategorical colimit universal property for groupoids now proves
the first assertion. The groupoid pushout of $1\leftarrow\mathbb Z\to1$
is the trivial groupoid, whereas the homotopy pushout of its classifying
spaces is the suspension of $S^1$. The pair $(\mathbb P^1,\infty)$ is the
geometric instance. This correction is consistent with F.13: its full
neighbourhood topoi have not been replaced by their covering reflections.
\end{proof}

\subsection{Finite flag truncations and van Kampen}

\paragraph{F.15 (Refuted). The two-depth diagram (246).}
For the coordinate divisor in $\mathbb C^3$, the bicategorical colimit
of the depth-one and depth-two groupoids in (246), with the incidence
isomorphisms but without the depth-three flag relations, has vertex
group $\mathbb Z^4$. The punctured normal boundary has fundamental
group $\mathbb Z^3$. These statements concern discrete groups; their
finite-covering analogues have the indicated profinite completions.

\begin{proof}
In polar coordinates the boundary of the real oriented blow-up is
\[
 L=(S^1)^3\times R,\qquad
 R=\{(r_1,r_2,r_3)\in\mathbb R_{\geq0}^3:\min r_i=0\}.
\]
The factor $R$ contracts to its vertex. Removing depth three deletes
that vertex, and radial normalization gives $R\setminus\{0\}\simeq
\partial\Delta^2\simeq S^1$. Thus $L_{\leq2}\simeq(S^1)^4$.
The index category has three singleton and three two-element subsets,
with six incidences: its nerve is a hexagon. Each total-link group is
$\mathbb Z^3$ and each incidence identifies its three coordinate
meridians. A spanning tree eliminates five incidence generators; the
remaining generator commutes with the three meridians, giving
$\mathbb Z^3\times\mathbb Z$. Adding the three-element subset and all
flags makes the index nerve a cone on that hexagon. Its six triangular
relations eliminate the six non-tree generators in the enlarged
presentation. This gives $\mathbb Z^3$.

The attached program constructs the subset posets, their edges and
triangles, chooses a spanning tree, and performs single-occurrence
Tietze eliminations; the outputs are respectively $(6,6,0)$ and
$(7,12,6)$ for vertices, edges, triangles, with one and zero surviving
index generators. This checks the actual bicategorical presentation,
not just its abelianization. The extra generator is precisely the
coherence lost by a strict colimit of one-object groups.
\end{proof}

\paragraph{F.16 (Proved). Sharp uniform depth bound.}
Let $L$ be the boundary of the real oriented blow-up of a smooth
normal-crossings pair, and let $L_{\leq r}$ be its part over points
with at most $r$ local branches. Include every incidence flag among
the retained depths. For $r\geq1$ the map
$L_{\leq r}\to L$ is $(r-1)$-connected, with the convention that an
$m$-connected map induces isomorphisms on $\pi_i$ for $i<m$ and a
surjection for $i=m$. Hence depths at most $k+2$ compute
$\tau_{\leq k}L$. This is the smallest uniform depth bound. The same
assertion holds on the classifying spaces of analytic
Deligne--Mumford stacks. In this last statement $L$ means the actual
oriented-blow-up boundary stack; identifying it with an intrinsic
punctured-neighbourhood germ still requires the hypotheses of F.19.

\begin{proof}
At a depth-$a$ point, the resolved boundary has $a$ angular coordinates
and a radial corner with at least one coordinate zero. Its depth-$a$
locus has real codimension $a-1$ in the boundary. Round corners and
triangulate compatibly. The deleted strata for $L_{\leq r}$ therefore
have codimension at least $r$. Relative general position moves a map
of a compact polyhedron of dimension less than $r$ off them, fixing
any closed subpolyhedron already disjoint. Apply this first to
representatives of homotopy groups and then to nullhomotopies. It
gives precisely the asserted connectivity. Local finiteness suffices:
the image of each such map meets only finitely many charts.

For sharpness take $m=r+1$ coordinate hyperplanes in $\mathbb C^m$.
The radial boundary is contractible, whereas its complement of the
vertex retracts onto $\partial\Delta^{m-1}$; consequently
\[
 L\simeq(S^1)^{r+1},\qquad
 L_{\leq r}\simeq(S^1)^{r+1}\times S^{r-1}.
\]
For $r=1$ the components differ. For $r>1$ the factor $S^{r-1}$
changes $\pi_{r-1}$. Thus $r\leq k+1$ cannot determine every $k$-type.
Exact integral boundary-matrix ranks for the simplex boundaries of
dimensions zero through four were computed independently; the radial
argument proves the formula in all dimensions. For a stack apply the
manifold result to every term of an etale atlas nerve. Realization
preserves an $m$-connectivity bound: relative CW models have cells only
above degree $m$, and realization adds nonnegative simplicial degrees.
This proves the assertion with isotropy retained.
\end{proof}

\paragraph{F.17 (Proved). Corollary 3
and the topoi that are truncated.}
Depths one, two and three, with every flag between them, compute
$\Pi_1(L)$. If $L\to X^*=X\setminus D$ is a one-type equivalence,
formula (247) has target $\Pi_1(X^*)$. The unstarred $\Pi_1(X)$ visibly
written in scan 152 is false under that hypothesis. For the full
\emph{sheaf pro-topos} gluing of F.13 all nonempty depths must be
retained; no dimension-independent cutoff computes it. The depth
bound in F.16 applies to the associated shape, not to that topos.

\begin{proof}
The fundamental-groupoid functor is the reflection from spaces to
one-types and preserves homotopy colimits, interpreted as
bicategorical colimits of groupoids. Apply F.16 with $k=1$. The
counterexample to the unstarred formula is $(\mathbb C,\{0\})$:
$L\simeq\mathbb C^*\simeq S^1$, whereas $\mathbb C$ is contractible.
This identifies a mathematical correction, not an alternative reading
of the ink. For full topoi a deleted deeper stratum is detected by
sheaves supported on it; even its locale of opens has changed. Thus
arbitrarily high depths cannot be removed in arbitrary dimension.
Ordinary sheaf descent uses objects, arrows and the cocycle condition
on triple overlaps; higher simplicial degrees are then forced. This
fact about \emph{simplicial coherence degree} does not permit omission
of high-depth objects: a depth-three object already supplies the
missing triangular relations in F.15.
\end{proof}

\paragraph{F.18 (Proved). Formula (254) in every dimension.}
With the germ conventions of F.12, put
$\mathcal V=\{\operatorname{Sh}(V)\}_{V\supset D}$ and
$\mathcal V^*=\{\operatorname{Sh}(V\setminus D)\}_{V\supset D}$.
Then in the pro-two-category of full Grothendieck topoi
\[
 \operatorname{Sh}(X)\simeq
 \operatorname{Sh}(X^*)\amalg_{\mathcal V^*}\mathcal V .
\]
For locally conical pairs the analogous gluing holds for sheaves of
spaces, and its shape is the homotopy pushout of the shapes of the
three terms. In particular
\[
 \Pi_1X\simeq\Pi_1X^*\amalg_{\Pi_1\mathcal V^*}\Pi_1\mathcal V.
\]
The two germ terms may be computed by the full flag diagrams of their
elementary pieces; for the punctured boundary one-type F.17 gives the
sharp depth-three replacement. For curves the closed term is the
disjoint union of the boundary points and the punctured term the
disjoint union of their meridian circles. This is the geometric form
of STRATA C.2 and C.5, not merely its presheaf special case.

\begin{proof}
For every neighbourhood $V$, $X^*$ and $V$ form an open cover of $X$.
Their sheaves are glued by identifying their restrictions to
$V\setminus D$, with the ordinary cocycle condition; F.13 verifies
this universally for geometric morphisms and passes to cofinal germs.
For space-valued sheaves use the entire Cech nerve: descent identifies
mapping spaces, not just their components. Applying this to a constant
space $K$ identifies maps from the shape with the homotopy limit of
maps from the three shapes. This is exactly the homotopy-pushout
universal property. Passing to one-types gives the displayed groupoid
formula. The tube retraction to a point and the punctured-disc
retraction to a circle give the curve case directly, without an
asphericity assumption on $X$. The intrinsic two-open argument also
applies to actual stack neighbourhoods. Computing those germs through
an atlas requires the proper or common-stage controlled class of F.13;
F.51 shows why the unrestricted atlas-germ substitution is invalid.
\end{proof}

\subsection{Normal directions, tangential points and arithmetic tubes}

\paragraph{F.19 (Proved). The canonical functor (243).}
Let $Y$ be a paracompact complex manifold with normal-crossings divisor
$D$, and let $D'\to Y$ be a normalized locally immersed intersection
of branches. Put $S=D'^*$, its open stratum, and let
$P_S=\prod_{j\in J,S}N_j^\times$ be the product of the punctured
normal line bundles. There is an equivalence
\[
 \Pi_1^{\mathrm{arcs}}(P_S)\longrightarrow
 \Pi_1^{\mathrm{rev}}(\mathcal V^*_{D',Y}),
\]
canonical as a functor with its specified natural comparison
isomorphisms. It is compatible with the generizations (231) and with
(234), including the coherence for every longer flag. No tubular
embedding is part of its definition. The intrinsic covering-germ
equivalence descends to the proper or common-stage controlled analytic
etale stacks of F.13. For an arbitrary etale presentation the sector
construction still defines a canonical coherent functor on its
chartwise germs, but these need not compute the intrinsic stack-open
neighbourhood germ. F.51 gives a parabolic quotient for which even
the latter covering equivalence fails. The quotient
$P_S/(\mathbb R_{>0})^J$ is the
open face over $S$ of the real oriented blow-up of $(Y,D)$.

\begin{proof}
Define the oriented blow-up intrinsically: a point over $y$ assigns
a unit complex argument to each local defining equation of a branch,
subject to $\arg(uf)=u(y)/|u(y)|\,\arg(f)$. A normal vector has the
same transformation rule. Its modulus is forgotten by the quotient
by positive real scalars; this gives the asserted face identification
without coordinates or a metric.

Here is the functor on coverings, which makes canonicity explicit.
Given a covering of a sufficiently small deleted neighbourhood and
a normal direction $v\in P_S$, choose a simply connected angular
sector about its arguments and a sufficiently small base neighbourhood
in $S$. The fibre at $v$ is the set of germs of sections on that
sector with all normal moduli positive and sufficiently small.
Shrinking sectors and radii does not change this set: in a normal
crossings chart these regions are contractible, and each covering is
a disjoint union of such regions. These sets form a covering of
$P_S$. A path in $P_S$ is covered by finitely many sectors; overlap
identification transports the germ of a section. Subdivision of the
path and homotopies give the same transport by uniqueness of lifting.
Thus evaluation gives a fibre functor on covering germs at every
$v$, and transport gives exactly (243).

In a chart $\Delta^a\times(\Delta^*)^{|J|}$ this operation identifies
both covering categories with $\mathbb Z^{|J|}$-sets, with the same
coordinate monodromies. It is fully faithful and essentially
surjective there. Coverings and their isomorphisms satisfy descent,
so these equivalences glue along all chart intersections. If $D'$
meets further branches, apply the same construction with their
coordinates also punctured. Pushing their positive radial coordinates
slightly into the interior shows that its covering-germ category is
unchanged upon replacing $D'$ by $S$; in each chart both categories
have exactly the same punctured coordinates and monodromy generators.
This proves (239) in the stated, not merely smooth-divisor, case.

For a nested flag use sectors with all its branch coordinates at once.
Successive restriction of a section germ is literal restriction, so
both orders agree; this proves (231), (234) and all longer
compatibilities before taking fundamental groupoids. On an etale
atlas all these constructions pull back and retain the isotropy
coherence. To identify the result with coverings of intrinsic stack
neighbourhoods, use the synchronized descent of F.13 in its stated
proper or controlled range. Without that hypothesis this argument
constructs the atlas-germ functor only; it does not silently exchange
an arbitrary infinite descent diagram with shrinking neighbourhoods.
In dimension one the construction evaluates a cover on a germ of a
sector tangent to a nonzero vector. The finite algebraic version is
Deligne's specialization to the punctured tangent line
\cite[\S15.13--15.27]{Deligne1989}. In a parameter $t$, it selects
compatible roots of $t$ and the fibre over the specified tangent
vector; changing $t$ by $u t$ multiplies the root coordinate by the
corresponding roots of $u(0)$. This is precisely the sector transition
rule. Hence the analytic, profinite, and tangential-point functors
coincide, with a path torsor rather than a falsely chosen canonical
basepoint identification. The specified cote~140-4 pages of
\emph{Longue Marche} were subsequently recovered inside the LM8
background archive; F.48 supplies the literal manuscript cross-check.
\end{proof}

\paragraph{F.20 (Refuted). An invariant tubular
embedding need not exist.}
The matrix in MPL~147, fol.~72, scan~146 must have upper-right entry
$1$ for the example to work. Let the discrete group $\mathbb Z$ act
on $\mathbb C^2$ by $(x,y)\mapsto(x+y,y)$ and let $D=\{y=0\}$.
There is no $\mathbb Z$-equivariant differentiable tubular embedding
from the normal line bundle to $\mathbb C^2$ which induces the
identity on its normal quotient. This is an example for an analytic
etale quotient stack with infinite isotropy, not for a proper
finite-stabilizer orbifold.

\begin{proof}
The action on $D$ and on its normal line is trivial. Differentiating
such an embedding at $(0,0)$ would give an equivariant splitting of
\[
 0\longrightarrow\mathbb C(1,0)\longrightarrow\mathbb C^2
       \longrightarrow\mathbb C\longrightarrow0.
\]
A fixed vector $(a,b)$ satisfies $(a+b,b)=(a,b)$, hence $b=0$;
none projects to $1$. This rules out even a real differentiable
splitting. The attached exact matrix calculation verifies this
obstruction. Entry $0$ would give the identity matrix and a splitting,
so it would not be a counterexample. For a finite group, averaging
an invariant metric supplies an equivariant smooth tubular
neighbourhood; the infinite stabilizer is material. The sector
construction in F.19 needs no such splitting and therefore survives.
\end{proof}

\paragraph{F.21 (Proved). The arithmetic
normal tube and (37).}
Let $(X,D)$ be a smooth normal-crossings pair over a characteristic-zero
field $K$, with $X$ of finite type, and let $S=D_J^*$ be a connected
open stratum on the normalization that labels its $r=|J|$ branches.
Define its tame etale tube by the finite-covering category of the
punctured henselian neighbourhood along $S$; equivalently use the
Kummer etale normal tube. This is independent of the chosen etale
charts. Its geometric profinite shape is that of
$P_S=\prod_{j\in J,S}N_j^\times$. With geometric basepoints and
paths specified, its fundamental groups have exact sequence
\[
 \pi_2^{\mathrm{et}}(S_{\bar K})
 \xrightarrow{\partial}
 T^J=\widehat{\mathbb Z}(1)^J
 \longrightarrow\pi_1^{\mathrm{et}}((\mathcal V_S^*)_{\bar K})
 \longrightarrow\pi_1^{\mathrm{et}}(S_{\bar K})\longrightarrow1.
\]
Here $\pi_2^{\mathrm{et}}$ is the profinite homotopy group of the
geometric etale shape, not the abstract completion of an unspecified
ordinary group. The map $\partial$ is the tuple of normal-line Chern
transgressions. Thus (37) is exact with $D_J$ meaning the open stratum;
its initial $1$ is valid exactly when $\partial=0$. Injectivity of
$T^J\to\pi_1^{\mathrm{et}}(U_{\bar K})$ implies that vanishing.
Arithmetic groups fit into the additional extensions over $G_K$;
Galois acts cyclotomically on each labelled factor of $T^J$ and
permutes factors if branch labels have not been split.

\begin{proof}
Locally along $S$ choose parameters $t_1,\ldots,t_r$ for its branches.
The henselian normal comparison sends a tame covering to its normal
specialization. A finite covering is trivialized locally by adjoining
roots $t_j^{1/n}$: normalization after these adjunctions is unramified,
by the valuation form of Abhyankar's lemma. On overlapping charts
$t'_j=u_jt_{\sigma(j)}$, the transition is the Kummer torsor of the
unit $u_j|_S$. These are exactly the transitions of $N_j^\times$.
This gives mutually inverse functors on finite covers and on their
isomorphisms, with the same descent cocycles; it proves the assertion
about the tube and its functor to $U$.

For the higher exact sequence, the root gerbes of the normal lines
provide a precise profinite model. For each $n$, the product of gerbes
of $n$th roots of the $N_j$ has fibre $B\mu_n^J$ over $S$ and is
classified by the Kummer Chern classes $c_1(N_j)\bmod n$. Their
compatible inverse system has fibre $B\widehat{\mathbb Z}(1)^J$.
The preceding local calculation, now for finite locally constant
truncated spaces by descent and their finite Postnikov towers,
identifies it with the profinite normal-tube shape. Equivalently this
is the local normal-bundle instance of the log comparison theorem
\cite[Theorem~A and Corollary~5.5, preprint numbering]{CSST2019}.
The homotopy exact sequence of this fibre sequence gives the displayed
sequence, and the boundary maps of the root gerbes are precisely the
Chern transgressions. Over $\mathbb C$ their evaluation on a sphere
is the integer degree of the pulled-back normal line bundle; reduction
modulo every $n$ gives the profinite statement. Descent from a finitely
generated field of definition and geometric extension of constants
preserves the finite-covering and finite-coefficient data, so the same
argument applies over arbitrary characteristic-zero $K$.

The composite $T^J\to\pi_1\mathcal V_S^*\to\pi_1U$ kills
$\operatorname{im}\partial$. If this composite is injective then
$\partial=0$, proving the requested direction. It need not vanish:
for the zero section in $\operatorname{Tot}\mathcal O_{\mathbb P^1}(1)$,
$\partial$ is an isomorphism and the punctured bundle is simply
connected; for a smooth conic in $\mathbb P^2$, its normal bundle has
degree four and the tube has fundamental group $\mathbb Z/4$, with
kernel $4\widehat{\mathbb Z}$ in its meridian group. Finally take
$X=\mathbb A^2$, $D=\{xy=0\}$ and the branch $D_1=\{x=0\}$.
The punctured tube has group $\widehat{\mathbb Z}^2$, its normal
meridian has rank one, and $D_1^*=\mathbb G_m$ supplies the other
factor. Replacing this open stratum by the closed $\mathbb A^1$
would make the last group trivial and break exactness. This also
confirms the mathematical force of (36) and the definition of $U_J$
as the trace of $U$ on the localized neighbourhood in scans~132--133.
\end{proof}

\paragraph{F.22 (Refuted). Finite roots do
not have a constant inertia kernel.}
Write $G=\pi_1^{\mathrm{tame}}(U_{\bar K})$ and
$\phi:T^J\to G$. The Kummer etale site of $(X,D)$ and the infinite
root stack have this same group and hence the same kernel
$\ker\phi$. For a finite root stack with orders $n_i$ along the
boundary components, put $N_{\mathbf n}$ equal to the closed normal
subgroup of $G$ generated by all $n_i$th powers of its corresponding
boundary inertia groups. Then
\[
 \pi_1(\sqrt[\mathbf n]{(X,D)})=G/N_{\mathbf n},\qquad
 \ker(T^J\to\pi_1\sqrt[\mathbf n]{(X,D)})
       =\phi^{-1}(N_{\mathbf n}).
\]
These kernels can be strictly larger than $\ker\phi$. Their
intersection over all divisible root orders is $\ker\phi$.
A finite Kummer \emph{cover} likewise cannot remove a nonzero kernel,
but this is a different assertion from equality for every finite
root stack.

\begin{proof}
A finite etale cover of the finite root stack is a tame cover of $U$
whose inertia orders divide the specified $n_i$. The assertion is
checked in a root chart $[\operatorname{Spec}A[z_i]/(z_i^{n_i}-t_i)
/\prod\mu_{n_i}]$: unramifiedness on the chart is precisely that
inertia bound; the charts and their covers descend. Therefore its
finite quotient actions are exactly those killing
$N_{\mathbf n}$, proving the displayed formulas. Every finite quotient
of $G$ kills $N_{\mathbf n}$ for some common multiple of its finite
inertia orders; thus $\bigcap_{\mathbf n}N_{\mathbf n}=1$.
All tame covers extend on the Kummer site by the same local
normalization argument, proving the infinite-level assertion.

For an explicit counterexample take $(\mathbb A^1,\{0\})$.
Here $G=T=\widehat{\mathbb Z}(1)$ and $\phi$ is the identity, while
$\sqrt[n]{(\mathbb A^1,0)}=[\mathbb A^1/\mu_n]$ has group $\mu_n$.
The kernel has changed from zero to $n\widehat{\mathbb Z}(1)$.
Finite quotients of this calculation were checked by the attached
program; the Kummer classification proves it in the inverse limit.

For a Kummer cover, the local lattice is a finite-index sublattice
$T'\hookrightarrow T$, and its global group embeds as an open subgroup
of $G$. Its kernel is the inverse image of $\ker\phi\cap T'$.
For any nonzero closed subgroup $A\subset\widehat{\mathbb Z}^{r}$,
and any finite-index sublattice $T'$, choose a positive integer $m$
with $mT\subset T'$. Multiplication by $m$ is injective on $T$, so
$0\ne mA\subset A\cap T'$. Thus the kernel persists. This proves
the valid content of EXIT\_ONE\_TYPE P1e-roots and refutes the stronger
finite-root-stack sentence added in the present prompt.
\end{proof}

\paragraph{F.23 (Proved). The stable cusp quotient changes the model.}
Over a characteristic-zero field containing $\mu_6$, put
$F=\{Y^2=X^3+1\}$ and
\[
 \mathcal X=[F\times\mathbb A^1_s/\mu_6],\qquad
 \zeta(X,Y,s)=(\zeta^{-2}X,\zeta^{-3}Y,\zeta s).
\]
Its boundary is $[F/\mu_6]$ at $s=0$; its open part is the family
$y^2=x^3+t$, $t\ne0$, via $t=s^6,x=s^2X,y=s^3Y$.
The normal meridian injects both discretely and profinitely into
the open fundamental group. Here groups are geometric; the discrete
calculation uses the complex realization of this model defined over
$\mathbb Q(\mu_6)$. The same holds for the full rank-two corner when
the infinity section in the projective elliptic completion is retained
as boundary. The original embedded cusp resolution has the exact
nonzero node kernels $\mathbb Z(-3,1)$ and $\mathbb Z(-2,1)$,
and their profinite counterparts, computed below. Thus the stable
quotient removes explicit local meridian kernels by changing the
model. It is not a finite root stack of that fixed resolved divisor.

\begin{proof}
The displayed substitutions verify the equation and the action is
free when $s\ne0$, so the asserted identification of the open part
is an isomorphism. A loop in the normal coordinate $s$ has degree six
under $t=s^6$. Its composite with the degree homomorphism to
$\pi_1(\mathbb G_m)=\mathbb Z$ is multiplication by six, which is
injective; the same calculation on $\widehat{\mathbb Z}(1)$ proves
profinite injectivity. No residual-finiteness inference from the
discrete injection is being used.

The affine curve $F$ is a once-punctured elliptic curve; its completion
is elliptic. Both are aspherical, and their finite quotient stacks
are aspherical by the fibration over $B\mu_6$. In the affine case
the coarse boundary coordinate $w=Y^2=X^3+1$ has cone orders two
and three, giving group $C_2*C_3$; the projective case adds the
order-six point at infinity. Thus the quotient retains finite
monodromy in boundary isotropy while retaining the infinite normal
meridian in the open part. F.22 shows why merely extracting finite
roots of the old boundary could not eliminate an existing kernel.
The construction uses inversion of $2$ and $3$; nothing here asserts
its validity for wild reduction at those primes.
\end{proof}

\paragraph{F.24 (Proved). The tame
constructible topos and its descriptive system.}
For a finite normal-crossings stratification over a characteristic-zero
field, finite etale sheaves constructible on its strata are finite
continuous representations of the profinite completion of its ordinary
exit category. The Ind-category of these finite representations is the
corresponding coherent constructible topos. Its stratum groups are the
geometric etale fundamental groups, with arithmetic descent over $G_K$;
its crossing sets and compositions are the profinite versions of the
coset and flag constructions in F.10. The groupoids of punctured
normal tubes are exactly the extensions in F.21, and the maps obtained
by forgetting branch flags are (38). These identify the claimed
``1-descriptif'' system without suppressing its kernels or its
commutation isomorphisms.

The arbitrary-product condition in F.8 is a characterization for
\emph{discrete} groupoid presheaf topoi. It is false as a necessary
condition for this profinite variant.

\begin{proof}
Over $\mathbb C$, finite constructible sheaves are recovered by their
finite stalks, path transport on each stratum, and the sector
specialization maps of F.19. The inverse construction glues constant
finite sheaves on the sectors; compatibility on double intersections
and on triple flags gives ordinary sheaf descent. Thus it is exactly
the set-valued case of F.10 with finite fibres. Every such finite
representation factors through a finite quotient of the category:
for each pair of the finitely many stratum basepoints, identify arrows
with the same map on all the specified finite fibres, and identify
endomorphisms with the same permutations. Composition descends because
these are a congruence. Taking all finite representations and their
finite products gives the cofiltered system of finite quotient
categories and proves the continuous-representation assertion.
The same construction for a finite family of representations proves
that the category is a pretopos, whose Ind-category is the asserted
coherent topos.

Analytification identifies the algebraic finite constructible sheaves
with these analytic ones, including their gluing functors, by
\cite[Proposition~12.6.4 and Corollary~12.6.7]{BGH2020}. More
explicitly, comparison holds on each smooth stratum for finite local
systems, and on its punctured normal charts by the Kummer calculation
of F.21; induction over the finite closed filtration identifies the
recollement triples, so it holds on the whole pair. Descent from a
field of definition and finite separable extension of constants gives
the arithmetic groupoids and their weak $G_K$-action as in F.1.
The descriptions of crossing sets and compositions follow by applying
finite stalk functors to the same sector and flag constructions.
In positive characteristic the same assertion concerns tame sheaves
and prime-to-characteristic local normal inertia; no wild cover is
silently included.

To see why products must be treated differently, let $G$ be a
profinite group and $H\subset G$ a closed subgroup which is not open.
In continuous discrete $G$-sets the product of a family is the subset
of the set-theoretic product whose points have open stabilizer. Choose
open subgroups $G_\alpha\supset H$ with intersection $H$. Each coset
$G_\alpha$ in $G/G_\alpha$ is fixed by $H$. Their tuple is not in the
continuous product, since its stabilizer is $H$. Therefore
\[
 \left(\prod_{\alpha}^{\mathrm{cont}}G/G_\alpha\right)^H
 \longrightarrow\prod_\alpha(G/G_\alpha)^H
\]
is not surjective. This is the actual cusp gluing functor for a
hyperbolic curve whenever its inertia subgroup has infinite index,
for example the thrice-punctured projective line. Nevertheless its
finite continuous exit description has just been proved. Finite
representations and coherent descent, rather than an unchanged
arbitrary-product test, give the valid profinite statement.
\end{proof}

\paragraph{F.25 (Refuted). Components of the analytic flag category
need not be a category.}
Even when the analytic incidence spaces satisfy all Segal identities,
their sets of connected components need not do so. The correct object
is the semisimplicial set of components of \emph{all} strict flags,
with repeated ranks adjoined if a simplicial object is desired.

\begin{proof}
Let $\tau=(12)(34)$ act freely on $\mathbb C^4\times\mathbb C^*$
by permuting the coordinates and sending $w$ to $-w$. Coordinate
hyperplanes descend to a normal-crossings divisor in a smooth
quotient. Every coordinate intersection times $\mathbb C^*$ is
connected, so connected flag components are the $\tau$-orbits of
branch-set flags. Put $B=\{1,2\}$. The upper pairs $(123,B)$ and
$(124,B)$ are in the same orbit, and the lower pairs $(B,1)$ and
$(B,2)$ are in the same orbit. However $(123,B,1)$ and $(123,B,2)$
are different orbits: applying $\tau$ changes $123$ to $124$.
Their composites $(123,1)$ and $(123,2)$ are also different orbits.
Thus a composable pair of component arrows has two possible
composites. The attached orbit enumeration checks these claims.
Forgetting entries of a complete flag always gives well-defined
face maps and satisfies the semisimplicial identities before and
after taking components; that is the stated replacement. Locally
its realization is the barycentric subdivision of the simplex of
branches, so its global realization is the usual coarse dual complex.
The normalized flag spaces themselves, with stabilizers and their
full fibre products, remain the correct geometric descent data.
\end{proof}

\subsection{Elementary topoi and the two shapes}

The intrinsic neighbourhood-germ statements in F.26--F.30 concern
analytic spaces and the proper or common-stage controlled stack
presentations specified in F.13 and F.19. The abstract exit-category
and bar-construction statements are separate from that geometric
identification. No assertion here replaces intrinsic germs by an
uncontrolled nonproper atlas, as ruled out by F.51.

\paragraph{F.26 (Proved). The four elementary pieces and their products.}
Over $(\mathbb C,\{0\})$ write
$P=\operatorname{Sh}(\{0\})$, $A=\operatorname{Sh}(\mathbb C^*)$,
$D=\{\operatorname{Sh}(\Delta_\epsilon)\}_\epsilon$ and
$D^*=\{\operatorname{Sh}(\Delta_\epsilon^*)\}_\epsilon$.
The canonical incidence maps \emph{over the ambient line} are
$P\to D$, $D^*\to D$, $D^*\to A$, their identities and composites
with the ambient inclusions. These generate exactly diagram (265).
There is no incidence $P\to A$ over the line. Structural maps to
the terminal topos exist, but they are not these ambient incidences.
For $I$ finite the elementary pieces are
\[
 E_{I_P,I_A,I_D,I_*}
 =P^{I_P}\times A^{I_A}\times D^{I_D}\times(D^*)^{I_*},
 \qquad I=I_P\amalg I_A\amalg I_D\amalg I_*.
\]
Their canonical incidence maps are the coordinatewise products of
the preceding three maps. For $|I|=2$ the complete list is
\[
\begin{array}{c|cccc}
 &P&A&D&D^*\\ \hline
 P&P P&P A&P D&P D^*\\
 A&A P&A A&A D&A D^*\\
 D&D P&D A&D D&D D^*\\
 D^*&D^*P&D^*A&D^*D&D^*D^*
\end{array}
\]
where juxtaposition denotes the product, not an identification.
The shape of such a piece is $(S^1)^{|I_A|+|I_*|}$, although its
full sheaf pro-topos is generally not the classifying topos of that
group.

Every canonical deleted germ (257), defined as the neighbourhood
germ of $\Sigma'$ in $\Sigma''$ followed by deletion of its trace
of $\Sigma$, is the bicategorical colimit of its contained elementary
pieces with their full incidence flags. This assertion uses full
sheaf pro-topoi. For shapes it is a homotopy colimit, and not the
shape of an arbitrary ordinary two-colimit of covering topoi.
For stack presentations this intrinsic-germ assertion has the proper
or common-stage controlled scope of F.13 and F.19.

\begin{proof}
The list is the independent choice of one of four states for each
coordinate: fixed zero, globally nonzero, near zero, and near zero
but nonzero. An inclusion over the ambient line can only fix zero
inside its germ, or include a punctured germ into that germ or the
nonzero line. This proves the generators and their coordinatewise
product description; it also explains why an apparent arrow from
$P$ to $\mathbb C^*$ in the typographical list (264) cannot be an
ambient inclusion. Polar retraction gives the displayed shapes.

At each positive radius the cover
$\mathbb C=\mathbb C^*\cup\Delta_\epsilon$ has intersection
$\Delta_\epsilon^*$. Its products form a finite cubical open cover,
and their intersections are exactly the stated punctured and
unpunctured coordinate pieces. Restriction to a coordinate closed
subspace replaces its zero coordinates by $P$. Restriction to an
admissible union, deletion of another admissible union, and taking
neighbourhoods of a third are computed in the same coordinates;
finite unions of tubes and their intersections give the corresponding
subdiagram. Apply full sheaf descent, including its triple-overlap
relations, and pass to the cofinal system of positive radius
functions. This is F.13 applied to this explicit cover, and proves
the assertion for (257). Applying space-valued descent gives the
homotopy-colimit assertion by the mapping-space argument of F.18.
On non-simple crossings in a Hausdorff analytic space use normalized
branch charts and their full etale nerve; their fibre products supply
the flag incidences and descend the same calculation. A stack atlas
additionally requires the synchronization hypothesis in F.13; no
unrestricted nonproper atlas-germ replacement is used. The order of taking a germ and deleting
its support matters: the punctured germ of a nonempty stratum is
not the germ of the empty stratum, exactly as scan~158 observes.
\end{proof}

\paragraph{F.27 (Proved). The Scholie and (295).}
For a general flag $d_*$ and two of its ranks $d'_0\geq d''_0$,
the elementary germ $\mathcal V^{d'_0}_{d_*,d''_0}$ has the
homotopy type of the punctured graded normal bundle over
$D_{d_*}^*$ with fibre $(\mathbb C^*)^{d'_0-d''_0}$.
Its fundamental groupoid is canonically equivalent, with the sector
comparison data, to that of this punctured bundle. This is (295).
Normal-line twisting and branch monodromy must be retained.

\begin{proof}
The two nested normal subbundles give an exact sequence
$0\to N'\to N''\to Q\to0$ as in (287). Remove the coordinate
hyperplanes from $Q$ and take its inverse image in $N''$.
The resulting map to $Q^\times$ is an affine bundle under the
pullback of $N'$; its fibres are contractible. Local trivializations
and a numerable cover show that this map is a homotopy equivalence,
without choosing a global splitting. Restrict to $D_{d_*}^*$.
The sector construction of F.19 compares its covering category with
the elementary germ and is compatible with refinement of the flag.
It therefore supplies the canonical equivalence on fundamental
groupoids claimed in (295). The homotopy equivalence need not be
a specified global analytic map, consistent with the warning at
(292). A line bundle of nonzero Chern class already shows why one
must keep the torus bundle rather than replace it by a product with
the base. For a stack the normal sequence and its punctured quotient
are formed on an atlas and descend; no invariant splitting is used.
That vector-bundle descent is distinct from identifying the result
with an intrinsic neighbourhood germ, for which the proper or
common-stage controlled hypotheses of F.19 remain necessary.
\end{proof}

\paragraph{F.28 (Proved). Exactly what the Remarque recovers.}
Let $\mathcal E=\operatorname{Exit}(X,D)$ be the full exit
infinity-category, with exits directed from deeper to more open
strata. The system of fundamental groupoids of the elementary
pieces, with every flag and its comparison isomorphisms, always
recovers $\Pi_1X$ and the topos of constructible \emph{sets},
$\operatorname{Fun}(h\mathcal E,\mathbf{Set})$.
If all open incidence strata are aspherical, the homotopy colimit
of the \emph{classifying spaces} of that same coherent system
recovers the homotopy type of $X$ and of every canonical germ.
Finally $\mathcal E$ is equivalent to an ordinary category if and
only if all these strata are aspherical and every full local
meridian lattice injects into $\pi_1(U)$. Under both hypotheses the
shape of the ordinary constructible topos is also the correct shape.
Asphericity alone does not imply this last conclusion.

\begin{proof}
The first groupoid assertion follows from F.18 and the preservation
of homotopy colimits by $\Pi_1$. For constructible sets, F.10 obtains
crossing Hom-sets as the components of the endpoint homotopy fibres;
only the fundamental groups of links and endpoints enter these
components. Full flag spaces give composition. This reconstructs
$h\mathcal E$ and its representation topos.

By F.27 every elementary link is a torus bundle over an incidence
cover of an open stratum, with possible additional contractible
factors. A finite cover of an aspherical space is aspherical, and
the long exact homotopy sequence of a torus bundle shows that its
total space is aspherical. Thus the natural maps from all diagram
values to their $B\Pi_1$ are equivalences, coherently with every
map and homotopy. F.26 now proves the second assertion.

For the criterion, endomorphism spaces in $\mathcal E$ are loop
spaces of strata; they are homotopy-discrete exactly when the strata
are aspherical. For incident $S,T$, choose its incidence cover
$F_{ST}\to S$ and normal link $L_{ST}$ with torus fibre. Normal
sectors give
\[
 \operatorname{Map}_{\mathcal E}(s,t)
 \simeq\operatorname{hofib}_{(s,t)}(L_{ST}\to S\times T).
\]
If the strata are aspherical, so is $L_{ST}$. A component of this
homotopy fibre is the classifying space of the kernel of
$\pi_1L_{ST}\to\pi_1S\times\pi_1T$. The incidence cover injects
on $\pi_1$, so this kernel is precisely that of the relative normal
lattice in $\pi_1T$. Hence all mapping spaces are zero-types exactly
when all relative meridian maps are injective.

It remains to justify that full injection suffices for every relative
map. Partition the local branches as $I=J\amalg K$, with the $J$
branches remaining zero on $T$. Push the full torus into the normal
torus bundle over $T$. If $v\in\mathbb Z^K$ dies in $\pi_1T$,
exactness lifts it to the image of some $w\in\mathbb Z^J$ in that
normal bundle. Then $v-w$ dies in $\pi_1U$. Full injection forces
$v-w=0$, and disjoint coordinate supports force $v=w=0$. This step
does not presuppose asphericity or injection of the intermediate
normal lattice. Taking $T=U$ proves the converse. It proves the
criterion, including all incidence covers and isotropy.
\end{proof}

\paragraph{F.29 (Proved). Shapes of the two constructible topoi and
the first obstruction.}
The genuine constructible infinity-topos has shape
$B\mathcal E\simeq X$; the infinity-topos associated with the
ordinary constructible-set topos has shape $B(h\mathcal E)$.
Their comparison
$q:X\to B(h\mathcal E)$ is two-connected. If every mapping-space
component is $(r-1)$-connected, $r\geq1$, it is $(r+1)$-connected.
For its homotopy fibre $F$, its first possible homotopy group is
\[
 \pi_{r+1}F\cong
 \operatorname{coker}\left[
 H_r(\widetilde V_2;\mathbb Z)
 \xrightarrow{d_0-d_1+d_2}H_r(\widetilde V_1;\mathbb Z)\right].
\]
Here $V_p$ is the bar construction of a well-based enriched model
of $\mathcal E$ with discrete objects, and tildes denote pullback of
the fundamental-groupoid local system of paths on the chosen
component of $B(h\mathcal E)$. The formula is equivariant for that
component's fundamental group.

\begin{proof}
In a presheaf infinity-topos on a category $C$, global sections of
the constant space $K$ are $\lim_CK=\operatorname{Map}(BC,K)$;
this proves the shape assertions. The exit category groupoidifies to
$X$ by conical sector descent, or by F.26 with its space-valued
coherence. Replace the exit category by a cofibrant enriched model
and put
\[
 V_p=\coprod_{x_0,\ldots,x_p}
 \prod_{i=1}^p\operatorname{Map}(x_{i-1},x_i),\qquad W_p=\pi_0V_p.
\]
The degree-zero map is an isomorphism. In positive simplicial degree,
$V_p\to W_p$ is $r$-connected. A relative cell first occurring there
has dimension at least $r+1$, and realization adds the positive
simplicial degree; the first relative cells thus have dimension
$r+2$. This proves $(r+1)$-connectivity, and $r=1$ always applies.

Pull back the universal path local system of $|W|$ to both bars;
faces transport the lift along the corresponding edge. Its relative
homology spectral sequence is
\[
 E^1_{p,j}=H_j(\widetilde W_p,\widetilde V_p;\mathbb Z)
 \Longrightarrow H_{p+j}(\widetilde{|W|},\widetilde{|V|};\mathbb Z).
\]
Column zero vanishes and all rows $j\leq r$ vanish. In the first
possible total degree $r+2$, the only position is $(1,r+1)$,
with $H_{r+1}(\widetilde W_p,\widetilde V_p)=H_r(\widetilde V_p)$.
Its incoming differential is the three-face map displayed above;
higher differentials have zero source or target. Relative Hurewicz
and the exact sequence of the homotopy fibre identify this cokernel
with $\pi_{r+1}F$. This supplies the topos version of
EXIT\_ONE\_TYPE P1d3 rather than merely referring to it.
\end{proof}

\paragraph{F.30 (Proved). The gluing question
and loss of homotopy.}
For a two-stratum conical space with closed stratum $S$, open stratum
$T$ and link $L$, the genuine constructible shape is
$S\amalg_L^hT$. When $S,T,L$ are aspherical with groups $A,B,G$,
the ordinary constructible shape replaces $G$ in this pushout by
$\operatorname{im}(G\to A\times B)$. Therefore the two shapes can
differ. For $(\mathbb P^1,\infty)$ they are respectively $S^2$ and
a point. For $(\mathbb P^n,\mathbb P^{n-1})$ the comparison fibre
is $\mathbb P^n$; for the square of $(\mathbb P^1,\infty)$ it is
$S^2\times S^2$. The first nonzero comparison-fibre group can occur
in arbitrarily high degrees, even for smooth algebraic pairs.

For a presheaf topos on $C$ with a full sieve $A$ and complementary
full subcategory $B$, let $E$ be the category of crossing arrows
$a\to b$, with commuting squares as morphisms. Then its shape is
\[
 BC\simeq BA\amalg_{BE}^{h}BB.
\]
This is the unconditional categorical answer to the question at
scan~180. For full geometric topoi the link is an actual topos if a
specified collared link exists; otherwise the deleted-neighbourhood
pro-topos is the precise object, as in F.12--F.18.

\begin{proof}
The conical mapping cylinder of $L\to S$ is a neighbourhood of $S$;
open descent with $T$ gives the first pushout. In the aspherical
case the crossing space is the homotopy fibre of
$BG\to B(A\times B)$. Its components are cosets of the image,
and their homotopy quotient by $A\times B$ is the classifying space
of the image stabilizer. Passing to the ordinary exit category
keeps only these components, which proves the replacement rule.
For the projective-line pair both endpoints are contractible and
the link is $S^1$. The ordinary category is $[1]$ and has
contractible nerve. For the hyperplane pair the open and closed
strata have trivial fundamental groups, and the Hopf-link endpoint
fibre is connected; again $h\mathcal E=[1]$. Products of exit
categories are computed coordinatewise, giving $[1]\times[1]$
for the square. These prove the three fibre assertions.

For arbitrary high degree take the smooth affine quadric
$Q_m=\{\sum_{i=0}^{m}z_i^2=1\}$, $m\geq2$, and the pair
$(Q_m\times\mathbb C,Q_m\times\{0\})$. Writing $z=x+iy$
identifies $Q_m$ with $TS^m$ up to the harmless positive rescaling
$|x|^2=1+|y|^2$, $x\cdot y=0$; scaling $y$ to zero retracts it
onto $S^m$. Its ordinary fundamental group is trivial.
The second factor's ordinary exit category has contractible
classifying space, so $B(h\mathcal E)$ is contractible while the
genuine shape is $S^m$. The first discrepancy is therefore degree
$m$, as asserted.

For the categorical statement, a nondegenerate string of arrows
in $C$ crosses from the sieve $A$ to $B$ at most once. Divide the
nerve into strings lying in either side and strings with their
unique crossing marked. The latter form the simplicial prism
resolution of $NE$; the two prism ends are its source and target
maps. Realizing this decomposition is the double mapping cylinder,
which proves the pushout, including disconnected and empty cases.
This checks STRATA C.3 and iterates to C.5. Its space-valued
coefficient form is the lax pullback of restriction functors:
objects are $(F_A,F_B,F_B\to i^*j_*F_A)$, and mapping spaces are
homotopy pullbacks. Restriction and the units of the two adjunctions
construct inverse functors, objectwise, proving STRATA C.6.
For locally constant coefficients the gluing arrow is an equivalence;
functors out of the just-proved homotopy pushout give the ordinary
infinity-categorical pullback of STRATA C.8. One cannot first take
homotopy categories: for the arrows $S^0\to *$ and $*\to S^1$,
the mapping-space components are
$\operatorname{coker}(\mathbb Z\xrightarrow{(1,1)}\mathbb Z^2)
=\mathbb Z$, whereas the ordinary comma category of homotopy
categories has a single morphism. This directly checks STRATA C.7
and locates the lost higher coherence.
\end{proof}

\subsection{Additional arithmetic tests and coherent limits}

\paragraph{F.35 (Proved). Inertia injection and extension as a sub-germ.}
Let $X/K$ be smooth in characteristic zero, $D$ a normal-crossings divisor,
$U=X\setminus D$, and $\bar x$ a geometric point of $D$. Choose a geometric
component of $U$ and a path from its base point to the strict local
punctured neighbourhood at $\bar x$. Write
\[
 H=\widehat{\mathbb Z}(1)^{I(\bar x)},\qquad
 G=\pi_1(U_{\overline K}),\qquad \iota:H\longrightarrow G.
\]
Then $\iota$ is injective if and only if every finite tame covering germ
at $\bar x$ is an open-and-closed sub-germ of the restriction of some
finite etale covering of $U_{\overline K}$. Equivalently, that covering
may be taken over $U$ after allowing a finite extension of the constants
and then composing its covering with $U_{K'}\to U$. The statement does
not assert that the entire restriction of the global covering is the
prescribed germ. This proves the first additional test F(h).

\begin{proof}
We prove the group-theoretic assertion for any continuous homomorphism
$\iota:H\to G$ of profinite groups. If it is injective, identify the compact
$H$ with a closed subgroup of $G$. A transitive finite $H$-set is $H/K$
for an open subgroup $K\subset H$. The compact subset $H\setminus K$
does not contain $1$, so some open normal $N\lhd G$ satisfies
$N\cap H\subset K$. Put $J=KN$. Since $N$ is normal, $J$ is a subgroup;
since it contains $N$, it is open. Also
\[
 J\cap H=K:
 \quad kn\in H\Longrightarrow n\in H\cap N\subset K.
\]
Consequently $hK\mapsto hJ$ embeds $H/K$ into the restriction of the
finite $G$-set $G/J$. Apply this to each of the finitely many orbits of
an arbitrary finite $H$-set. Conversely, if $1\ne h\in\ker\iota$, some
finite continuous quotient of $H$ detects $h$. Its regular $H$-set
cannot embed in a restricted $G$-set, on which $h$ acts trivially.
This proves the equivalence.

Finite tame germs are finite continuous $H$-sets by the strict local
Kummer calculation of F.21. Finite etale coverings of the geometric
component of $U$ are finite $G$-sets. An invariant subset of a finite
covering is precisely an open-and-closed subcover, so the group assertion
has exactly the stated geometric interpretation. The global covering and its finite presentation descend to a finite
extension $K'/K$; in characteristic
zero this extension is separable. Composition with the finite etale map
$U_{K'}\to U$ gives the stated version over $K$. The comparison remains
geometric: for a nonclosed point one must not substitute the whole
arithmetic decomposition group for $H$.
\end{proof}

\paragraph{F.36 (Proved). The coherent etale limit is nondegenerate.}
Let $x$ be a closed point on a smooth curve over an algebraically closed
field $k$ of characteristic zero. Let $R=\mathcal O_{X,x}^{\mathrm{sh}}$,
$F=\operatorname{Frac}R$, and let $(V_i,\bar x_i)$ run through pointed
affine etale neighbourhoods of $\bar x$, shrunk so that a local parameter
$t$ is defined and its only zero in $V_i$ lies over $x$. Put
$V_i^*=V_i[1/t]$. Then
\[
 \varprojlim_i (V_i^*)_{\mathrm{et}}
   \simeq (\operatorname{Spec}F)_{\mathrm{et}}
   \simeq B\widehat{\mathbb Z}(1)
\]
as a bicategorical limit of Grothendieck topoi. The hypotheses of the
coherent-limit assertion in the LM8 intake are therefore sufficient.
The exact general references are SGA~4, VI, Theorems~8.2.3 and~8.3.13,
and VII, Lemma~5.6; the last statement is also the Stacks Project,
Lemma~59.51.2 \cite{SGA4II,StacksEtaleLimits}.

\begin{proof}
The schemes $V_i^*$ are affine, the transition maps are affine, and
$\varinjlim_i\mathcal O(V_i)=R$. Localization commutes with that colimit,
so their scheme limit is $\operatorname{Spec}R[1/t]=\operatorname{Spec}F$.
An affine etale algebra of finite presentation over this limit is given
by finitely many coefficients and finitely many equations; it descends
to some $V_i^*$. The same finite-presentation argument descends morphisms
and equalities. A covering can be replaced by a finite affine covering
of each affine object, and surjectivity of that finite family descends
after increasing $i$. Thus the restricted etale site of the limit is
the filtered colimit of the restricted sites, including the covering
families. Passing to sheaves gives exactly the inverse-topos universal
property of VI~8.2.3.

Each stage is a coherent topos and every transition morphism is coherent.
One can see the extra assertion of VI~8.3.13 directly here: a coherent
object is the quotient of a finite diagram of coherent generators by a
finitely presented equivalence relation. Generators, relation, and its
finite reflexivity, symmetry and transitivity identities all descend to
one stage. Morphisms descend by the same finite presentation. Hence
coherent objects in the limit are the filtered two-colimit of coherent
objects at the stages; no analytic infinite annulus covering is present.

The ring $R$ is a strictly henselian discrete valuation ring. Its residue
field has characteristic zero, so all finite extensions of $F$ are tame.
Every unit of $R$ has an $n$th root by Hensel's lemma. The tame extension
classification therefore gives the separable closure
$\bigcup_n F(t^{1/n})$, with Galois group
$\varprojlim_n\mu_n=\widehat{\mathbb Z}(1)$; the classification follows
by removing the ramification index with $t^{1/n}$, after which the
extension is unramified and hence trivial over a strictly henselian
ring. This also follows from the tame valuation calculation in the
Stacks Project, Lemma~15.116.2 \cite{StacksTameDVR}.
The topos is not degenerate: its terminal finite set is not its empty
set. In residue characteristic $p$, the tame quotient is instead
$\prod_{\ell\ne p}\mathbb Z_\ell(1)$ and the full group has wild inertia.

There is no conflict with F.12 or LM8 E.11. The classical shrinking-disc
system has the annular covering that kills every cone. The etale system
has finite coherent covering descent and the nonempty scheme limit
$\operatorname{Spec}F$. MPL~147, scans~143--148 and 157--169, explicitly
works with analytic neighbourhoods, arcs, and analytic multiplicities;
its $D,D^*$ are the classical pro-objects. The phrase ``tubulures etales''
of MPL~149, scans~133--135, proposes this different arithmetic construction.
\end{proof}

\paragraph{F.37 (Refuted). Pinned fullness over a finite field.}
Even after localization at universal homeomorphisms, the etale-topos
functor on finite-type $\mathbb F_q$-schemes is not full onto pinned
morphisms over $B\mathbb F_q$. This does not contradict F.4, which proves
faithfulness. It also explains why the infinite-field hypothesis in
\cite[Theorem~0.1]{CHW2024} must not be silently dropped.

\begin{proof}
Let $X=\mathbb A^1_{\mathbb F_q}$, let $p:X_{\mathrm{et}}\to
B\mathbb F_q$ be the structure morphism, and identify
$\Gamma_{\mathbb F_q}$ with $\widehat{\mathbb Z}$. Since
$\pi_1(X)\to\Gamma_{\mathbb F_q}$ is onto and the latter is abelian,
every $a\in\widehat{\mathbb Z}$ gives a natural automorphism of $p^*$:
on a finite constant-field covering it acts by $a$. Thus
\[
 (\operatorname{id}_{X_{\mathrm{et}}},a):
 X_{\mathrm{et}}\longrightarrow X_{\mathrm{et}}
 \quad\text{over }B\mathbb F_q
\]
is a relative geometric morphism. Its underlying map of spaces is the
identity, so it is pinned.

We claim that the identity geometric morphism of $X_{\mathrm{et}}$ has
no nontrivial natural automorphism. Such an automorphism restricts to
every nonempty open of $X$, and then to its generic point. Every finite
separable extension of $\mathbb F_q(t)$ spreads to a finite etale cover
of some such open, so restriction gives an element of
$Z(G_{\mathbb F_q(t)})$. This restriction is injective. Indeed affine
etale schemes are generators; an automorphism of such a representable
which is generically the identity is the identity, because its source
is reduced and its affine target is separated. Affine generators are
essential here: one need not make a generic-detection assertion for
arbitrary nonseparated etale objects.

To prove $Z(G_{\mathbb F_q(t)})=1$, let a central element $z$ project to
$\sigma\in\Gamma_{\mathbb F_q}$. Write
$N=G_{\overline{\mathbb F}_q(t)}$. For any prime $\ell\ne p$, the Kummer
classes of $t-b$, $b\in\overline{\mathbb F}_q$, are linearly independent
in $\overline{\mathbb F}_q(t)^*/(\overline{\mathbb F}_q(t)^*)^\ell$,
as their valuations at the distinct points $b$ show. The conjugation
outer action of $\sigma$ on $N$ permutes these classes by
$b\mapsto\sigma(b)$, with the scalar factor prescribed by the action
on $\mu_\ell$. If $\sigma$ moves any $b$, this cannot be the identity
on $\operatorname{Hom}(N,\mu_\ell)$: distinct valuation basis vectors
cannot become one another by a nonzero scalar. Centrality therefore
forces $\sigma=1$. Hence $z\in Z(N)$, which is trivial by the valuation
and Kummer proof of F.2, unchanged in characteristic $p$ when its auxiliary
prime is chosen different from $p$. We have proved the claim.

An isomorphism between the two relative maps indexed by $a,b$ would now
have to be the identity automorphism of the underlying identity map;
compatibility with their base isomorphisms forces $a=b$. There are
therefore uncountably many isomorphism classes of pinned relative maps
of this form. In contrast, the category of finite-type schemes over a
finite field has a countable skeleton and countable Hom-sets: all its
affine presentations and finite gluing data use finite strings over a
countable set. Its localization at universal homeomorphisms still has
countable Hom-sets, since morphisms are represented by finite zigzags.
It cannot surject onto this uncountable set of pinned morphisms.
This is a complete cardinality counterexample; it does not require
choosing an unspecified nonintegral $a$ and asserting that particular
one is nonalgebraic. Powers of Frobenius account for some, but not all,
of the relative twists.
\end{proof}

\paragraph{F.38 (Known). The verified reconstruction range over larger bases.}
The following are the precise positive ranges used for MPL~149,
scan~52. Over an infinite finitely generated field, use pinned morphisms
and invert universal homeomorphisms (F.5). The same holds over every
sub-$p$-adic field, meaning a subfield of a finitely generated extension
of $\mathbb Q_p$. More generally, let $S$ be noetherian, Nagata and
geometrically unibranch, with sub-$p$-adic function field on every
irreducible component. On finite-type $S$-schemes which are
\emph{componentwise universally open}, the localization at universal
homeomorphisms is fully faithful into topoi over $S_{\mathrm{et}}$ with
\emph{chivalrous} morphisms: the induced continuous map sends
constructible subsets to constructible subsets. In particular normal
flat finite-type schemes over $\operatorname{Spec}\mathbb Z$ are
reconstructed, with all their morphisms. These are Theorem~4.1,
Proposition~6.11 and Theorem~6.12 of the opened September~20, 2026 version
of Carlson--Stix \cite{CarlsonStix2026}; the relative hypothesis is part
of the theorem and is not a theorem about all schemes over $S$.

\begin{proof}
We spell out the reduction from the quoted statements to the formulation
at scan~52. A universal homeomorphism induces an equivalence of small
etale topoi, so the functor factors through the stated localization.
Proposition~6.11 identifies that localization, in the componentwise
universally open category, with absolute weak normalizations. Theorem~6.12
identifies the Hom-set from $X^{\mathrm{awn}}$ to $Y$ with the groupoid
of chivalrous relative topoi morphisms; its rigidity assertion says that
this groupoid is equivalent to that set, with no extra automorphisms.
For a normal $X$ in this category, $X=X^{\mathrm{awn}}$ by their
Lemma~6.10. Thus the asserted Hom-equivalence follows, and applying it
to a topos equivalence and its inverse gives inverse scheme morphisms
by faithfulness. Flat finite-type schemes over $\mathbb Z$ satisfy the
componentwise openness condition, giving the stated example.

Over a field all these schemes are componentwise universally open and
finite-type field morphisms send closed points to closed points. Their
Theorem~4.1, together with the seminormalization description of the
localization in characteristic zero, gives the pinned field assertion.
In the general-base theorem, pinned is the wrong substitute for
chivalrous: for a discrete valuation ring $R$ and nonunit $a\ne0$, the
closed point $(aT-1)$ of $\mathbb A^1_R$ maps to the generic point of
$\operatorname{Spec}R$. This genuine scheme morphism is chivalrous by
constructibility of images, but not pinned. The proof of the relative
theorem reconstructs over the generic fields and spreads out using
normal quasi-sections; this is exactly where geometric unibranchness,
Nagata finiteness and componentwise openness enter. Thus its application
here does not discard any of those hypotheses.

Combining F.5--F.7 and F.37 gives the necessary qualifications to the
provisional question: unrestricted fullness is false even over number
fields; pinned fullness holds in the preceding ranges, fails over finite
fields, and fails over algebraically closed characteristic-zero fields.
Faithfulness of genuine scheme maps remains the all-characteristic
statement F.4. The remaining general-field question is recorded precisely
in F.44, not identified with the false finite-field assertion.
\end{proof}


\paragraph{F.39 (Proved). Stable towers, full inertia and cyclotomic normalization.}
Let $U\to B$ be an elementary hyperbolic fibration over a field $K$
finitely generated over $\mathbb Q$, with $B$ a connected smooth
finite-type $K$-variety: a smooth proper curve family with
a nonempty finite etale marked divisor removed, of type $(g,n)$ with
$2g-2+n>0$. Then each horizontal geometric puncture group injects into
$\pi_1(U_{\bar K})$. For a recursively compatible tower of these
fibrations, equip its base with a normal-crossings completion whose
strata are aspherical and whose full meridian lattices inject both
discretely and profinitely. Suppose the fibration extends to a stable
pointed family there, after finite root or quotient changes, with a
smooth toroidal total model and fibre incidence type constant on each
base stratum. Then every total full meridian group injects:
\[
 I_{\bar x}=\widehat{\mathbb Z}(1)^{I(\bar x)}
       \lhook\joinrel\longrightarrow\pi_1(U_{\bar K}).
\]
These assertions are natural for etale changes and compatible monomial
maps, with path conjugacy retained. On a splitting field of the
branches, conjugation on each normal factor is the cyclotomic character;
without splitting, it is the product of that action with the branch
permutation action. A monomial map $t'_j=u_j\prod_i t_i^{m_{ji}}$ induces
the integer contact matrix $(m_{ji})$ on inertia. This supplies a common
geometric injection and normalization theorem for the arithmetic issue
raised in the sixth EXIT task. F.48 audits the recovered
\texttt{CYCLOTOMIC\_KERNEL} rows; the remaining literal comparison
with the unlocated \texttt{FLAGGED\_149} rows is recorded in F.45.

\begin{proof}
We give the completion argument, since discrete injection alone is
insufficient. Every stage of the tower has this smooth finite-type
base geometry. After choosing $\bar K\hookrightarrow\mathbb C$, the
proper smooth family, with its finite etale markings, is locally a
trivial differentiable family of marked compact surfaces: a horizontal
distribution tangent to the markings is obtained by a partition of
unity and its local flows give the trivializations. Removing the
markings therefore gives a locally trivial analytic topological bundle.
Its fibre is a punctured surface with group a free
group $N=F_r$, $r\ge2$. For a bundle over any connected base, the image
of $\pi_2(B)\to\pi_1(F)$ is central: pull back over a sphere, write its
clutching loop as $h_t\in\operatorname{Homeo}(F)$, and use the square
$h_t(\gamma(s))$ to commute the evaluation loop $h_t(x)$ with every
$\gamma\in\pi_1(F,x)$. The free group's centre is trivial, so the fibre
injects. A puncture loop is nontrivial: if there are several punctures
it survives in homology; with one puncture it is the reduced nonempty
product of handle commutators. Thus each such loop has infinite order.

We prove that for \emph{every} extension
$1\to F_r\to G\to Q\to1$ the completed sequence is exact on the left.
First $Z(\widehat F_r)=1$. Given a finite quotient $F_r\to A$, with
kernel $L$, choose a prime $\ell\nmid|A|$. The cellular chain complex
of the finite regular covering of the $r$-rose gives
\[
 0\longrightarrow H_1(L,\mathbb F_\ell)
 \longrightarrow\mathbb F_\ell[A]^r
 \longrightarrow\mathbb F_\ell[A]
 \longrightarrow\mathbb F_\ell\longrightarrow0.
\]
Averaging over $A$ splits these representations, so
$H_1(L,\mathbb F_\ell)\simeq
\mathbb F_\ell\oplus\mathbb F_\ell[A]^{r-1}$ is a faithful $A$-module.
The finite quotient $F_r/[L,L]L^\ell$ consequently detects the image in
$A$ of any purported central element of $\widehat F_r$: centrality
forces its action on this module to be trivial, and hence its image
in $A$ is $1$. This for every $A$ proves centrelessness.

The automorphism group of a finitely generated profinite group is
profinite: the intersections of its finitely many open subgroups of
index at most $m$ form a characteristic open basis, and automorphisms
are the compatible invertible maps on the resulting finite quotients.
The conjugation map
$G\to\operatorname{Aut}(\widehat F_r)$ therefore extends to
$\widehat G$. Its composite with
$\widehat F_r\to\widehat G$ is the inner-automorphism map, which is
injective by centrelessness. Thus $\widehat F_r\to\widehat G$ is
injective. Its image is the closure of $F_r$, and the quotient of
$\widehat G$ by that closure is $\widehat Q$, since both classify the
finite quotients killing $F_r$.

A nontrivial cyclic subgroup $\langle w\rangle\subset F_r$ has its full
profinite topology. Indeed represent a cyclically reduced conjugate
of $w$ by an immersed labelled circle in a rose. For every label, pair
the missing incoming and outgoing edges to complete this finite graph
to a finite covering of the rose, without identifying its vertices.
The original circle is an embedded cycle in that graph, so its integral
homology class is primitive. The corresponding finite-index subgroup
$H\subset F_r$ has a homomorphism $H\to\mathbb Z$ sending $w$ to $1$.
For any $m$ the core in $F_r$ of the kernel of $H\to\mathbb Z/m$
gives a finite quotient in which $w$ has order divisible by $m$.
This proves $\widehat{\langle w\rangle}\simeq\widehat{\mathbb Z}
\hookrightarrow\widehat F_r$. Together with the extension argument it
proves horizontal profinite injection over an arbitrary base. By
induction a tower is aspherical, and its discrete group is residually
finite: separate an element either in the completed base or in the
embedded completed free fibre.

For the vertical assertion let $A$ be a total local normal lattice,
$A_B$ its base lattice, and $M:A\to A_B$ the integral monomial map.
At a smooth vertical point $\ker M=0$; at a marking it is generated
by the fibre puncture; at a node $xy=t$ it is generated by $(1,-1)$
in the map $(u,v)\mapsto u+v$. In the latter case its image in the smooth
fibre is the vanishing cycle. Root changes and smooth toroidal charts
replace these lattices by finite-index or primitive sublattices; hence
in general $\ker M$ has rank at most one and maps to a nonzero power
of the corresponding cycle. The cycle is essential. If a neck bounded
an unpunctured disc, the side it cut off would be an unmarked genus-zero
tree attached at one node, impossible for a stable pointed curve:
the sum of the special-point counts on such a tree is less than three
times its number of components. A neck around one marked point remains
a nontrivial puncture loop, and the extra two-punctured bridges of a
toroidal model also have essential meridians.

The diagram
\[
\begin{array}{ccccc}
 \ker M&\longrightarrow&A&\longrightarrow&A_B\\
 \downarrow&&\downarrow&&\downarrow\\
 \pi_1(F)&\longrightarrow&\pi_1(U)&\longrightarrow&\pi_1(B)
\end{array}
\]
has the needed exactness in its rows. An element of $A$ killed in
$\pi_1(U)$ maps to zero in the injective base lattice, then belongs to
$\ker M$, where the essential-cycle map is injective. It is zero.
Complete the lower sequence by the free-kernel result. Smith normal
form proves $\ker\widehat M=\widehat{\ker M}$: the image of $M$ is
free abelian, the sequence onto it splits, and its inclusion into
$A_B$ remains injective after completion. The same diagram chase,
using completed base injection and cyclic fibre injection, proves the
full profinite assertion. The strata are bundles of punctured stable
components or tori over finite incidence covers of base strata; their
homotopy sequences prove their asphericity. These statements give the
induction for every recursively compatible tower in the theorem.

The Riemann existence comparison is SGA~1, XII, Theorem~5.1 and
Corollary~5.2 \cite{SGA1}. It identifies these completed groups and
Kummer meridians with the geometric etale groups, so it transports
the proved injections. In local parameters the automorphisms of
$t_i^{1/n}$ are multiplication by $\mu_n$; Galois conjugation acts on
these roots of unity by the cyclotomic character. Taking all $n$
gives the asserted Tate modules. Pulling back $t'_j$ gives the contact
matrix, while the roots of the units $u_j$ alter the tangential
splittings by their Kummer cocycles. They do not alter the homomorphism
on geometric meridians. Thus neither a preferred generator of
$\widehat{\mathbb Z}(1)$ nor a preferred tangential path has been
asserted.

Finally, if a local arithmetic extension
$1\to I_x\to E_x\to G_{\kappa(x)}\to1$ maps to that of $U/K$ and the
chosen map $G_{\kappa(x)}\to G_K$ is injective, then
\[
 \ker(E_x\to\pi_1(U))=
 \ker(I_x\to\pi_1(U_{\bar K})).
\]
Project a kernel element to $G_{\kappa(x)}$ to prove this equality.
The condition holds at a closed point after its finite residue-field
extension, but is not automatic at a transcendental point.

There are two precise limits on this theorem. First, cyclotomic action
on a subgroup is a necessary normalization of genuine inertia; we have
not asserted that it characterizes inertia among all subgroups.
Second, stable extension alone is insufficient for an arbitrarily
chosen base completion. Blow up a smooth boundary point of the product
completion of $(\mathbb P^1\setminus\{0,1,\infty\})^2$. At the new
double point the exceptional and strict-transform meridians both map
to the old meridian, so $(1,-1)$ is killed. A constant stable curve
family over that base retains this kernel, discretely and profinitely.
The compatibility and full base-lattice assumptions above are therefore
necessary to this induction. This also verifies, rather than merely
quotes, the distinction between EXIT P3b-vertical and its counterexample.
\end{proof}

\paragraph{F.43 (Proved). Finite-homotopy exodromy for $\overline M_{0,5}$.}
Let $X=\overline M_{0,5,\mathbb C}$ with its normal-crossings boundary
stratification, and let $C=h\operatorname{Exit}(X^{\mathrm{an}})$.
The continuous functors from the strict profinite completion of $C$ to
$\pi$-finite spaces classify the etale sheaves of $\pi$-finite spaces on
$X$ constructible for these strata. Stalks and all generization maps are
preserved. Thus the $n=5$ question of LM3b H.9 and Part E(c) is decided;
one must retain the category and its flags, not just its groupoidification.

\begin{proof}
Use
\[
 U=M_{0,5}=\{(x,y):x,y\in\mathbb P^1\setminus\{0,1,\infty\},\ x\ne y\}.
\]
The compactification is the blow-up of $\mathbb P^1\times\mathbb P^1$
at $(0,0),(1,1),(\infty,\infty)$. Forgetting $y$ is a locally trivial
bundle with base the thrice-punctured sphere and fibre the
four-times-punctured sphere: small disjoint discs around the moving
puncture give local trivializations by an isotopy supported in those
discs. The homotopy exact sequence proves that $U$ is aspherical and
that
$1\to F_3\to\pi_1(U)\to F_2\to1$ is exact. All ten one-dimensional
boundary strata are thrice-punctured spheres and the fifteen deepest
strata are points, so every stratum is aspherical.

Let $a,b,c,d,e$ denote the homology classes of meridians of
$x=0,x=1,y=0,y=1,x=y$. They are a basis of $H_1(U,\mathbb Z)$: project
the five complex affine line complements to their defining linear
functions to see independence, and attach the normal disc boundaries
to the simply connected affine plane to see generation. The ten
boundary vectors are
\[
\begin{array}{c|c@{\qquad}c|c}
x_0&a&y_0&c\\
x_1&b&y_1&d\\
x_\infty&-a-b-e&y_\infty&-c-d-e\\
\Delta&e&E_0&a+c+e\\
E_1&b+d+e&E_\infty&-a-b-c-d-e.
\end{array}
\]
The infinity formulas are the product-one meridian relations on a
projective line; an exceptional meridian is the sum of the three
meridians through the blown-up point. The intersecting pairs are
$(x_i,y_j)$ for $i\ne j$, and $(E_i,x_i),(E_i,y_i),(E_i,\Delta)$.
For each of these fifteen pairs the $5\times2$ matrix has a $2\times2$
minor $1$ or $-1$ after an integral change of rows; equivalently the gcd
of its maximal minors is $1$. The attached exact computation lists all
fifteen gcds and obtains Smith invariants $(1,1)$ in each case.
Consequently every local meridian lattice $\mathbb Z^r$, $r\le2$,
embeds as a direct summand of $H_1(U,\mathbb Z)$, hence injects into
$\pi_1(U)$ and, after completion, into $\widehat{\pi_1(U)}$.

For any incident strata $S,T$, let $L$ be the normal link with its two
maps. The exit mapping space is the homotopy fibre of
$L\to S\times T$. Its higher homotopy groups vanish because these
three spaces are aspherical. Its fundamental group is
$\ker(\pi_1L\to\pi_1S\times\pi_1T)$; an element in that kernel is a
normal meridian and is zero by the injection just established (with the
relative-face argument of F.28). Thus every component of every exit
mapping space is contractible. The exit infinity-category is therefore
equivalent to its ordinary category $C$, with composition specified by
the actual three-stratum flags. This is the extra step not furnished by
asphericity of strata alone.

We next justify completion; it is not a formal assertion that profinite
completion preserves all homotopy pullbacks. A finitely generated free
group is good: the cellular two-term resolution of a rose computes its
cohomology with every finite module, and the same free resolution in
continuous modules computes the cohomology of the free profinite group.
The full profinite topology on $F_3$ inside $\pi_1(U)$ is induced by
conjugation into $\operatorname{Aut}(\widehat F_3)$, as proved in F.39.
Thus the profinite completion of the displayed free-by-free sequence is
exact. Comparing its discrete and continuous Hochschild--Serre spectral
sequences proves goodness of $\pi_1(U)$: the cohomology groups of the
fibre with finite coefficients are finite, and only rows $0,1$ occur,
so goodness of the free quotient applies to every term. Links over a
one-dimensional stratum have group $\mathbb Z^r\rtimes F_2$ (the
extension splits since $F_2$ is free), and links over a point have group
$\mathbb Z^r$. Their completions are exact by embedding the normal
lattice into $\widehat{\mathbb Z}^r\rtimes\widehat F_2$; the finite
Koszul resolution of $\mathbb Z^r$ and the same spectral-sequence
comparison prove goodness. Hence all the stratum and link profinite
homotopy types are the classifying types of their profinite groups.
The completed meridian injection now proves, by the same homotopy-fibre
calculation, that every profinite exit mapping object is a profinite
set. Descent on the finitely many full flag diagrams reconstructs a
profinite ordinary EI category with these Hom spaces and compositions.

Here is the finite-congruence point explicitly. A compact totally
disconnected category with finitely many objects is the inverse limit
of its finite categorical quotients. Indeed, choose finite clopen
partitions of its Hom spaces, and put $a\sim b$ when every defined
left and right context $u(-)v$ puts $uav$ and $ubv$ in the same partition
piece. Joint continuity and compactness of the context spaces imply
that this equivalence is open: finitely many neighbourhoods of contexts
suffice to find one neighbourhood of $a$ on which all their partition
values are fixed. Compactness then makes the number of classes finite.
Inserting further contexts shows that the equivalence is compatible
with composition. Thus finite congruences are cofinal among the finite
clopen data. In our case their restrictions to the dense discrete
category $C$ are precisely the strict finite congruences used in LM3
F.2--F.4. Here one also uses the topology assertion in F.39 and the
completed link descriptions above: finite quotient data on the
discrete stratum groups and incident cosets extend continuously to
their completions, and the finitely many composition constraints
extend by density. Density alone would not imply this extension
property. Thus the inverse system is the strict profinite completion.

The same statement works with finite homotopy values, not only finite
sets. Given finitely many $\pi$-finite values, take a degreewise finite
simplicial model of the full infinity-category they span; its mapping
spaces have one common finite truncation bound. Its coherent diagram
data are determined by finitely many simplicial degrees, using a
coskeletal model above that bound. Continuous maps in those degrees
factor through finite clopen partitions of finite products of Hom
spaces. Every such partition admits a finite coordinatewise clopen
refinement: compactness gives a finite cover by clopen rectangles,
and the finitely many possible object sequences cause no additional
infinite choice. Refine the resulting individual Hom partitions by
the context congruence of the preceding paragraph. Thus one finite
categorical congruence suffices for all the diagram data.
The same argument with an additional finite simplex
$\Delta^m$ proves the assertion on mapping spaces, for every $m$.
Consequently continuous representations in $\pi$-finite spaces are the
infinity-categorical filtered colimit of their finite-stage
representations; this is stronger than a statement about finite sets.

Finally, the algebraic-to-analytic comparison on these strata and their
normal links identifies their finite covers and their finite-coefficient
higher descent data. In the notation of Barwick--Glasman--Haine, this
is Proposition~12.6.4 and Corollary~12.6.7, followed by the exodromy
Theorem~10.1.8 \cite{BGH2020}. The preceding paragraphs verify in this
particular geometry that the resulting profinite exit object is the
strict completion of the ordinary category, rather than a profinite
infinity-category with hidden higher mapping spaces. Applying the
comparison and exodromy therefore gives precisely the asserted
stalk-preserving equivalence. The classifying space of this category
still has higher homotopy; for example $H^2(X,\mathbb Z)=\mathbb Z^5$
and $H^4(X,\mathbb Z)=\mathbb Z$. Ordinary exit mapping spaces being
discrete does not mean that this classifying space is a $1$-type.
\end{proof}

\paragraph{F.44 (Open). Delimitation beyond the proved reconstruction ranges.}
The requested base-field delimitation now includes the following
proved positive ranges: infinite finitely generated fields by F.5,
sub-$p$-adic fields and the specified general bases by F.38, and
every algebraic subfield of $F\mathbb Q^{\mathrm{ab}}$, for a
number field $F$, by F.47, and, more generally, every subfield of
a finitely generated extension of $\mathbb Q^{\mathrm{ab}}$ by
F.50. These assertions include arbitrary cyclotomic extensions of
number fields and fields finitely generated over them. In particular, the test
over $\mathbb Q(\mu_{2^\infty})$ considered during this work is
settled, including full site realization. F.46 records its initial
point-function proof and F.47 supplies the additional argument
needed for the full geometric morphism. The proposed unrestricted
fullness is false by F.5. F.49 strengthens the counterexamples
F.6 and F.37: pinned fullness fails over every algebraic extension
of a finite field, every algebraically closed field, and every
infinite perfect field with nontrivial centre of its absolute
Galois group, including every real closed field.
Faithfulness of actual reduced-scheme maps remains
the all-characteristic theorem F.4.

The exact general criterion is Carlson--Haine--Wolf,
Theorem~4.18 \cite{CHW2024}: a field has pinned reconstruction if
and only if every pinned morphism from a connected regular affine
scheme with a rational point to $(\mathbb G_m)_{\mathrm{et}}$
has the geometric-point values of a regular unit (on the absolute
weak normalization in positive characteristic). Its implication to
full site realization is proved through all sliced etale objects,
as explained in F.47. Their Theorem~6.24 further supplies the
positive range of amenable fields admitting a nontrivial discrete
valuation. Amenability means that rational points of every torus
have no nonidentity prime-to-characteristic divisible element and
that the prime-to-characteristic Tate module of the Picard group of
every regular finite-type scheme vanishes (Definition~5.15).
This is an exact verified sufficient condition, not a claim that
every field outside it is a failure or a literature-wide open case.

A concrete test not decided by this return is
\[
 k_{\mathrm{solv}}=\mathbb Q^{\mathrm{solv}},\qquad
 \operatorname{Hom}_{k_{\mathrm{solv}}}(\mathbb G_m,\mathbb G_m)
 \longrightarrow
 \pi_0\operatorname{Hom}^{\mathrm{pin}}_{B G_{k_{\mathrm{solv}}}}
       ((\mathbb G_m)_{\mathrm{et}},(\mathbb G_m)_{\mathrm{et}}),
 \tag{F44}
\]
where $\mathbb Q^{\mathrm{solv}}$ is the maximal solvable algebraic
extension in a fixed algebraic closure. The map is injective by
F.4; its surjectivity is not asserted. Finite-type dimension-zero
schemes, after universal-homeomorphism localization, are finite
etale and their relative Galois topoi already determine their maps.
Thus (F44) is a first positive-dimensional test of the remaining
criterion, rather than a claim of a smallest base field or of
literature-wide minimality. Constant morphisms and the general
nonconstant, possibly noninvertible, locus are part of the question;
the unresolved realization issue concerns the latter locus.

Here are the precise failed extensions of the proof, rather than
an appeal to the absence of a cited theorem. The field
$k_{\mathrm{solv}}$ is radically closed. Indeed, an element lies
in a finite solvable Galois extension $E/\mathbb Q$. Adjoin the
needed roots of unity and then $n$th roots of all its conjugates;
the latter extension is abelian over the enlarged $E$, so its
normal closure is still solvable over $\mathbb Q$. Consequently
\[
       D(k_{\mathrm{solv}}^*)=k_{\mathrm{solv}}^*.
\]
The height argument used in F.47 cannot reduce these elements to
roots of unity: this field contains $2^{1/n}$ for every $n$, of
positive heights tending to zero. The same observation shows that
it cannot be contained in $F\mathbb Q^{\mathrm{ab}}$ for any
fixed number field $F$, by the verified Amoroso--Zannier bound.
It is also outside the broader range F.50: the algebraic constants
in a finitely generated extension of $\mathbb Q^{\mathrm{ab}}$
form a finite extension of that field, to which the same height
bound applies.
It is not amenable. It is not sub-$p$-adic either, since it contains
all roots of unity whereas the algebraic constants in a finitely
generated extension of $\mathbb Q_p$ form a finite extension and
have only finitely many roots of unity.

The Picard input of F.47 also fails here. For an elliptic curve
$E/\mathbb Q$, every $2^r$-division field is solvable: its Galois
group is a subgroup of
$\operatorname{GL}_2(\mathbb Z/2^r)$, an extension of
$\operatorname{GL}_2(\mathbb F_2)\simeq S_3$ by a $2$-group.
Hence $E[2^\infty]\subset E(k_{\mathrm{solv}})$, providing
nonzero infinitely divisible points and a nonzero Tate module.
Neither the integral-residue argument followed by Picard
divisibility nor the translated-function proof can therefore be
copied with its hypotheses unchanged. Kummer classes at the
$k_{\mathrm{solv}}$-points already forget all nonzero constants.

The relative birational isomorphism results checked in
Tsujimura's Theorem~4.7 and Corollary~4.10 \cite{Tsujimura2022}
do not supply this noninvertible realization step; the requisite
AVKF hypotheses also fail in this example. Descent of a scheme to
a number field does not imply descent of an arbitrary geometric
morphism on the whole etale site. No exotic pinned morphism in
(F44) has been constructed here, and no argument using the
nonabelian finite covers has yet eliminated one. The next step
would have to exploit those nonsolvable-cover compatibilities,
rather than the Kummer and Picard invariants just shown to vanish
or to have the wrong divisibility. This is the exact residual
test and obstruction left by this return; no broader assertion
about the status of all other fields is made.

A further literature check does not settle (F44).
Mochizuki's \emph{Global Solvably Closed Anabelian Geometry}
\cite{Mochizuki2006SolvablyClosed}, Definitions~1.1 and~2.3 and
Theorem~2.5, concerns quotient groups
$Q_F=\operatorname{Gal}(\widetilde F/F)$ for a solvably closed
extension of a number field, together with the required
$\Sigma$-compatible quotients of hyperbolic-curve fundamental
groups. Its conclusion reconstructs compatible isomorphisms.
The absolute group $G_{\mathbb Q^{\mathrm{solv}}}$ in (F44)
is instead a kernel in $G_{\mathbb Q}$, and (F44) asks about
all pinned morphisms, including noninvertible ones. Neither
changing the quotient to this kernel nor replacing the Isom
conclusion by Hom is justified by that theorem.

\medskip
\noindent\emph{Finite extensions do not repair these particular inputs.}
One possible approach would pass to a proper finite extension
$L/k_{\mathrm{solv}}$, which is Hilbertian by the Weissauer theorem
used explicitly in the proof of \cite[Proposition~2.5]{F44Jarden2011}.
Descent of a reconstruction theorem would indeed be valid here.
More generally, in characteristic zero pinned reconstruction
descends along a finite extension $L/k$. To prove this, apply the
criterion of \cite[Theorem~4.18]{CHW2024} to a connected regular
affine $X/k$ with a rational point and a pinned morphism to
$(\mathbb G_m)_{\mathrm{et}}$. The curve-free criterion applies
after finite etale base change to $L$, where it gives a unit
$f_L$ with the prescribed geometric-point values. On
$X_{L\otimes_kL}$ the two pullbacks of $f_L$ have the same values
at every geometric point, by the descent data of the original
morphism. This scheme is reduced, so the units agree.
Faithfully flat descent gives $f\in\Gamma(X,\mathcal O_X)^*$
with the required values over $k$; the same criterion proves
reconstruction. The source remains connected after base change:
regularity in characteristic zero makes it normal and smooth,
and its rational point forces its finite algebraic constant field
to be $k$.

The obstacle is therefore the arithmetic input over $L$, not this
descent argument. For \emph{every algebraic} extension
$L/k_{\mathrm{solv}}$, including each proper finite one, one has
\[
 k_{\mathrm{solv}}^*\subset D(L^*),\qquad
 E[2^\infty]\subset E(L).
\]
Thus $L^*/\mu_\infty$ contains the nonzero divisible subgroup
$k_{\mathrm{solv}}^*/\mu_\infty$ and cannot be free abelian;
the same elliptic curve still violates the AVKF and Tate-module
conditions. Nor can such an $L$ admit a nontrivial discrete
valuation. A homomorphism $v:L^*\to\mathbb Z$ arising from a
valuation vanishes on the divisible subgroup $k_{\mathrm{solv}}^*$.
For any element algebraic over this trivially valued field, its
minimal polynomial has all nonzero coefficients of value zero.
A positive or negative value of the element would make the
minimum value among its nonzero polynomial terms occur exactly
once, which is impossible. Hence $v$ vanishes on $L^*$.
Hilbertianity of a finite extension therefore does not restore
the multiplicative, Picard, or discrete-valuation hypotheses used
in the reconstruction proofs verified here. This eliminates that
specific route without producing an exotic pinned morphism or
deciding (F44).


\medskip
\noindent\emph{Additional audit of the residual test
$k=\mathbb Q^{\mathrm{solv}}$.}
Neither of the elementary nongeometric constructions in F.49 settles
this test. In fact, with $G=G_{\mathbb Q}$ and
$H=G_{\mathbb Q^{\mathrm{solv}}}$, one has
\[
                         C_G(H)=1.                         \tag{F44.C1}
\]
We give the proof, verifying the precise scope of the recovered
\texttt{CYCLOTOMIC\_KERNEL} row B2 instead of importing a stronger
unqualified assertion about arbitrary closed subgroups.

The following argument works for every Hilbertian field $F$ and
every nontrivial closed normal subgroup $N\triangleleft G_F$.
Two primary results suffice. Jarden's Proposition~2.5 says that if
$E/F$ is algebraic Galois and $E$ is not separably closed, then
$G_E$ is almost locally free; in particular it contains a free
profinite subgroup of rank two
\cite{F44Jarden2011}. Haran--Jarden, Corollary~2.5, says that
$F^{\mathrm{sep}}$ cannot be the compositum of two proper subfields
which are Galois over $F$ \cite{F44HaranJarden1991}. Both statements
allow infinite Galois extensions.

There is no nontrivial closed abelian normal subgroup $A$ of $G_F$:
its fixed field is Galois over $F$ and is not separably closed, so
the first result would place a nonabelian free profinite rank-two
group inside $A$. Put $C=C_{G_F}(N)$. It is closed and normal,
and $N\cap C=Z(N)$ is a closed abelian normal subgroup of $G_F$;
hence $N\cap C=1$. If $C\ne1$, the two proper fixed fields of
$N$ and $C$ would be Galois over $F$, and their compositum would
have absolute Galois group $N\cap C=1$. This contradicts the
second result. Thus $C_{G_F}(N)=1$.

To apply it to $F=\mathbb Q$, we check that
$\mathbb Q^{\mathrm{solv}}\ne\overline{\mathbb Q}$.
The polynomial $P=X^5-X-1$ is irreducible modulo $3$: the attached
exact program divides it by all twelve monic polynomials of degree
one or two and finds no divisor. A reducible quintic must have a
factor of one of these degrees. Modulo $2$ its squarefree
factorization is
\[
           (X^2+X+1)(X^3+X^2+1),
\]
with both factors irreducible. Its rational Galois group is therefore
transitive of degree five and contains a permutation of cycle type
$(2)(3)$, whose cube is a transposition. A transitive group of prime
degree is primitive. The graph whose edges are the conjugates of
that transposition has components forming blocks, and is therefore
connected. Its edge transpositions generate $S_5$, so the Galois
group is $S_5$. The program also computes its derived-subgroup
orders $120,60,60$, certifying nonsolvability. Its splitting field
cannot lie in the maximal solvable extension. Thus $H\ne1$, and
Hilbert irreducibility for $\mathbb Q$ gives (F44.C1).
In particular $Z(H)=1$, so F.49 supplies no nonidentity central
comparison twist over this field.

The semilinear field-automorphism construction also does not apply.
Since $k/\mathbb Q$ is Galois,
$\operatorname{Aut}_{\mathbb Q}(k)=G/H$. For an automorphism
$\sigma$ choose a lift $g\in G$. Its outer action on $H$ is
conjugation by $g$. If this action were inner, say equal to
conjugation by $h\in H$, then
\[
                       h^{-1}g\in C_G(H)=1.
\]
Hence $g=h$ and $\sigma=1$. A nonidentity coefficient-semilinear
automorphism of $\mathbb A^1_k$ consequently cannot be supplied with
the relative comparison used in F.49: such a comparison, restricted
along the rational point $0$, would make its base automorphism of
$BH$ two-isomorphic to the identity and thus its action on $H$
inner. This excludes that construction; it makes no assertion that
all pinned topoi morphisms over $k$ have been classified.

The finite certificate is
\texttt{code/F44\_solvable\_certificates.py}; its exact output is
included with the deliverable.


\paragraph{F.45 (Open). The remaining literal FLAGGED\_149 audit.}
The required literal comparison with the \texttt{FLAGGED\_149}
register rows F11--F16 and F25 remains incomplete. No exact matching
return was found in the supplied files, the targeted title and content
searches, or the complete 749-item inventory of accessible ZIP
filenames. That inventory does not exclude an unindexed copy nested
inside another archive. The actual register and its proofs are
therefore still needed to verify the intake's row numbers and claims.

Two earlier source gaps were closed during this work. The
\texttt{LM8\_LIE\_APPROXIMATIONS\_background.zip} archive contains
the actual \texttt{CYCLOTOMIC\_KERNEL} return and the \S53 scans and
transcriptions. F.48 audits its exact rows B5 and C8 and proves the
literal sector-functor comparison with \S53. Earlier claims that
these particular sources had not been located are superseded by
that recovery. The initial unsuccessful searches are retained only
as a record of the work, not as current descriptions of availability.

The EXIT rows P3b-discrete, P3b-completion,
P3b-horizontal-profinite, P3b-vertical,
P3b-vertical-counterexample and P3b-profinite were inspected with
their proofs; F.39 gives the relevant stable-tower theorem and its
precise contact-matrix and cyclotomic qualifications. F.22 refutes
the separate assertion that finite root stacks have one fixed
uncompleted inertia kernel. These independent mathematical results
do not authenticate the unlocated FLAGGED register. Only that
literal cross-return comparison remains open in this source-audit
row; F.48 closes the previously recorded CYCLOTOMIC and \S53 gaps.
